% Français 9115, batch 21 — Gallica views 401–420.
% Pass: Opus 5.5 (claude-opus-5-5), 2026-09-22 — first pass, unchecked against the views by a human.
%
% What the twenty views hold. Nineteen views carry text and one is blank.
% Unlike the recto-only runs of batches 13 to 16, these leaves are written
% on both sides: the verso of each ink-numbered leaf carries the next page
% of her pagination. There is no repeated exposure and no microfilm
% target. With one interpolated slip (v. 416) and one sheet of rough
% calculation (v. 410–411), the whole batch is a single piece in her fair
% hand, lightly corrected: a memoir on the vibrations of elastic
% surfaces, paginated (4) to (17) at the head, which begins before this
% batch (view 400, just outside it, is page (3), inked 196, and computes
% A, B, C for the first of the two planes) and goes on after it (view 421
% is page (18)).
%
%   v. 401–404  End of her § 1: the angle e between two consecutive
%               planes through four neighbouring points of the surface,
%               sin² e = (dx²+dy²)(d²z)²/ds², on neglecting dz for a
%               surface near a plane; the moment ∫∫Eβ δe, with the force
%               of elasticity taken proportional to the sum of the inverse
%               radii of curvature, E = bτ³{1/(r)+1/(r')}, b the natural
%               rigidity and τ the thickness; integration by parts,
%               boundary terms set aside; the equation of motion (A),
%               bτ²{d⁴z/dx⁴ + 2d⁴z/dx²dy² + d⁴z/dy⁴} + d²z/dt² = 0 (v. 403);
%               N.o 2, the regular motion likened to a pendulum of length
%               K, giving (B), Kbτ²{…} = z; with Kbτ² = h⁴, (B) splits into
%               two second-order equations, one of them that of tensed
%               surfaces, so that every nodal figure possible for a tensed
%               surface is possible for an elastic one within the same
%               bounds (v. 404).
%   v. 405–409  § 2, « Intégration de l'équation differentielle de la
%               surface elastique vibrante ». N.o 3: the time factor
%               sin(t√(1/K)+ζ) understood thereafter; first integral (C)
%               for d²z/dxdy = 0 (exponentials, sines and cosines of
%               ωx/A, ωy/A), its sounds √(1/K)/π = ω²τ√b/(πA²). N.o 4:
%               periodic surfaces, d²z/dx² = (n²/m²)d²z/dy², integral (D),
%               A sin(πnx/A) sin(πmy/A), sounds π(m²+n²)τ√b/A²; then the
%               series integrals (E), (F), in powers of cos, with a
%               parameter Π, and a new condition between d²z/dx² and
%               d²z/dy². N.o 5 (v. 407–409): z = z'z'', the integral (G)
%               under Π + Π' = Π'', (H) with coefficients 2−m², 2−n²; N.o 6
%               (v. 409): z' periodic in x only, integral (I); all these
%               integrals but the first also belong to tensed surfaces,
%               whence their simplicity.
%   v. 410–411  Not part of the fair copy: one leaf of rough calculation,
%               v. 410 written across the sheet (photographed sideways),
%               v. 411 its verso, differentiating a series of the form of
%               (F) twice in x. No pagination; inked 201 on v. 410.
%   v. 412–415  Fair copy resumed, and a second copy of what v. 408–409
%               hold: (G) again, (H) again, then (I), N.o 6 and the remark
%               on the simplicity of these integrals (pages (11), (12));
%               page (13) ends it — the other integrals would be harder to
%               apply — and opens § 3, « Observations sur les propriétés
%               analytiques des differens points des corps sonores », N.o 7
%               (v. 414): a line satisfying the analytic conditions of the
%               contour may be taken as a limit of a vibrating part; for the
%               string, the midpoints of the ventral segments satisfy them
%               as well as the nodes, since R.(dz/dx).δz = 0 is met by
%               z = 0 or dz/dx = 0 (« formule P 332 de la 1ere édition de la
%               Mec. Analy. »), confirmed for wind instruments by
%               « M.r Dela Grange » in « le second volume des mémoires de
%               Turin »; the open flute is the reverse of the string (v. 415).
%   v. 416      A small torn slip, pasted on a guard, written in a finer,
%               paler pen: an error « dans tous le passage relatif aux
%               effets de l'application dela cire » on the sound of a plate
%               loaded with wax, (N²+n²)/n² against (N²+n²)/N², and 2n²/n²
%               against 2N²/(N²+n²), « qui heureusement ne different pas
%               sensiblement ». It interrupts the memoir: v. 415 stops on
%               « la manière » and v. 418 goes on « de vibrer la plus
%               simple… ».
%   v. 417      Blank verso of the slip of v. 416 (checked mirrored: only
%               the slip's writing, reversed). Not transcribed.
%   v. 418–420  Pages (15)–(17). Tensed surfaces: « M.r Biot P. 87 deson
%               memoire t. 4. del'institut »; not every line satisfying the
%               extremity conditions is a line of rest; for tensed surfaces
%               the nappes and lines of rest answer to Sauveur's « ondes »
%               and « nœuds » « ainsi que l'indique le programme publié par
%               la classe », but lamina and elastic surfaces have two kinds
%               of lines of rest; a remark « qui n'a pas encore été faite »,
%               arising from Euler's theory of the lamina, « Acta academiæ …
%               Petropolitanæ pro anno 1779. Pars prior » (v. 419). N.o 8
%               (v. 420): after Lagrange's extremity conditions (« P. 197 de
%               la nouvelle éd. de sa mec. », 197 doubtful), four couples of
%               equations; Euler (« P. 116 du mémoire cité ») examined only
%               three; the fourth, d³z/dx³ = 0 and dz/dx = 0, also gives
%               free ends, and « je crois être parvenue à obtenir ces
%               derniers. mais la lame […] s'est brisée trop tôt pour que je
%               pusse constater et mesurer ce résultat ». The leaf ends
%               there; view 421, page (18), continues.
%
% The numbers on the leaves, and why no \folio{} is written. (1) At the
% top right of every ink-numbered recto, in the larger, rounder hand and
% broader pen of the running series: 197 (v. 402), 198 (404), 199 (406),
% 200 (408), 201 (410, read with the sheet turned), 202 (412), 203 (414),
% 204 (416), 205 (418), 206 (420) — one a leaf, no gap, continuing 196 on
% view 400. The versos (401, 403, 405, 407, 409, 411, 413, 415, 417, 419)
% carry none. Following the house rule of this volume, these penned
% numbers are recorded in notes and here, never as \folio{}. (2) Her own
% pagination, in parentheses at the head: (4) v. 401, (5) 402, (6) 403,
% (7) 404, (8) 405, (9) 406, (10) 407, (9) 408, (10) 409, (11) struck 412,
% (12) 413, (13) struck 414, (14) 415, (15) 418, (16) 419, (17) 420. Views
% 410, 411 and 416 carry none. The (9) and (10) are read twice (406/407
% and 408/409), and (11)–(12) repeat the matter of 408–409: v. 408–409
% and 412–413 are two successive fair copies of the second half of (10)
% onward. Recorded as they stand; not reconciled. (3) Her article numbers
% « N.o 2 » to « N.o 8 » and section heads « §. 2. », « §. 3. » are text.
% No pencilled foliation was seen on any view. Every number above was read
% on its view; none is inferred.
%
% Laubenbacher and Pengelley's « Manuscript C » (Fermat for even exponents,
% ff. 348r–349r) is not in this batch: nothing here is number theory.
%
% Notation. Hers, kept. Squares and cubes of differentials written with the
% figure above the letter (dx², d²z) are set as ordinary exponents. The
% variation sign, a d with a looped stem, is set δ (hand.md, v. 140); her
% double integral, two long s, as ∫∫. The thickness is a stroke with a
% rightward cap, like Γ or τ, set τ (doubtful at its first occurrence,
% v. 402). In the series of v. 407–413 one tall-legged π serves both for
% the half-circumference in the arguments and for a parameter in the
% coefficients; it is set π in the first use and Π in the second (Π', Π''),
% a distinction of the transcriber's. A looped capital coefficient is set
% \mathcal{A} (v. 405–406) and an ornate one \mathfrak{A} (v. 409–410, 413).
% The point between two figures is her product sign (3.5 = 15, 3.8 = 24).
% « sin », « cos » set as \text{sin} etc., stops kept where written. The
% large brace and bracket layouts of the series (E)–(I) are simplified into
% aligned lines; the terms are given in her order.
%
% Spelling is hers and is not flagged: « demême », « dela », « delà »,
% « vis avis », « mouvemens », « differens », « fesant », « quarrés »,
% « s'en suivra », « c'est adire », « ayent », « déja », « longtems »,
% « aproximation », « aprecier », « prens », « parvenue » (feminine),
% « cequi » for « ce que », « donné » for « donnée », « tous le passage »,
% « limites donnés », and « Dela Grange ». She writes Euler's E like an L (« Lobr », v. 419 — the
% stroke hand.md records at v. 210–218).
%
% What could not be read. \ill{} stands in five places: a letter struck
% after « que » in the phrase « dont jouit l'integrale » (v. 404); a
% struck group in the numerator of the last condition (v. 407); the
% blackened factor after « dz/dx = » (v. 411); the looped-over group after
% « Π + Π' = » (v. 412); the struck first spelling under « exponentielles »
% (v. 413). Readings offered with doubt are marked \uncertain{}: τ (v. 402),
% « lasuite » (v. 406), « may » (v. 409), « aient » and « 24 » (v. 412),
% « mn », « faudroit », « sur » (v. 413), « de » (v. 415), « cordes » and a
% struck « s » (v. 419), « 197 » and « mesurer » (v. 420). Where the leaf and
% the calculation disagree the leaf stands, with a note (v. 407 and 412,
% arguments in mx where my is due; v. 410, (Π+3)(Π+3.4) in one line).
%
% Edges. View 401 continues page (3) on view 400 (« on a demême pour le
% second plan »); view 420 is continued by page (18) on view 421.
%
% Nothing here was checked against any published edition of Sophie
% Germain, Euler, Lagrange or Biot.
\documentclass[11pt,a4paper]{article}
% The preamble shared by every transcription of the mathematicians' manuscripts in Gallica.
%
% It rests on one idea: **uncertainty compounds, it does not resolve.** A
% transcription that smooths over a doubtful word has destroyed the only thing
% it was made for — a reader must be able to tell, at every point, what was on
% the page and what was supplied. The apparatus macros below are therefore the
% critical apparatus, not decoration.
%
% The same commands are recognised by `scripts/render.mjs`, which produces the
% site's reading view, and by `scripts/tei.mjs`. A macro added here must be
% added there too, or rendering fails loudly — which is the wanted behaviour.
%
% Comments here are English, like the rest of the repository. The documents
% this preamble sets are French, and that is deliberate: see the skills.

\ProvidesPackage{germain}[2026/09/15 Transcriptions of mathematicians' manuscripts in Gallica]

\usepackage{amsmath,amssymb,amsthm}
\usepackage{mathrsfs}
% Kept from the parent preamble so the renderer's diagram support stays
% available; these manuscripts draw no commutative diagrams, and nothing here uses it.
\usepackage{tikz-cd}
\usepackage[english,french]{babel}
\usepackage{fontspec}

% --- Greek ------------------------------------------------------------------
% The text face stays fontspec's default, Latin Modern, which has no Greek:
% XeTeX drops every glyph it lacks with a line in the log and nothing on the
% page, and the Greek words the manuscripts quote vanished from the PDFs. So
% the Greek and Greek Extended blocks get a XeTeX character class of their own,
% and entering or leaving a run of it switches the family to CMU Serif — the
% same Computer Modern design, with polytonic Greek — and back. Only the
% family changes: a Greek word in \textit or \textbf keeps its shape and series.
% The transitions to and from every other class carry the tokens that class
% has towards an ordinary letter, so French spacing before « ; » or inside
% « » still holds after a Greek word. No grouping, so no brace or \endgroup
% can come out unbalanced; maths is left alone, since XeTeX inserts nothing
% there. Kept identical in `exercices.sty`.
\makeatletter
\newfontfamily\germainGreekfont{cmun}[
  Extension=.otf, NFSSFamily=germaingreek,
  UprightFont=*rm, ItalicFont=*ti, SlantedFont=*sl,
  BoldFont=*bx, BoldItalicFont=*bi, BoldSlantedFont=*bl]
\newXeTeXintercharclass\germain@greek
\def\germain@greekrange#1#2{%
  \@tempcnta=#1\relax
  \loop
    \XeTeXcharclass\@tempcnta=\germain@greek
    \ifnum\@tempcnta<#2\advance\@tempcnta\@ne\repeat}
\germain@greekrange{"0370}{"03FF}
\germain@greekrange{"1F00}{"1FFF}
\def\germain@greekprev{\rmdefault}
\def\germain@greekin{%
  \edef\germain@greekprev{\f@family}\fontfamily{germaingreek}\selectfont}
\def\germain@greekout{\fontfamily{\germain@greekprev}\selectfont}
\def\germain@greekjoin{%
  \XeTeXinterchartokenstate=\@ne
  \@tempcnta=\z@
  \loop
    \ifnum\@tempcnta=\germain@greek\else
      \XeTeXinterchartoks\germain@greek\@tempcnta=\expandafter{%
        \expandafter\germain@greekout\the\XeTeXinterchartoks\z@\@tempcnta}%
      \XeTeXinterchartoks\@tempcnta\germain@greek=\expandafter{%
        \the\XeTeXinterchartoks\@tempcnta\z@\germain@greekin}%
    \fi
    \ifnum\@tempcnta<4095\advance\@tempcnta\@ne\repeat}
% babel-french sets its own classes, and saves and restores the interchar
% state, each time a language is selected: join again after it does.
\addto\extrasfrench{\germain@greekjoin}
\addto\extrasenglish{\germain@greekjoin}
\AtBeginDocument{\germain@greekjoin}
\makeatother

\usepackage{geometry}
\geometry{a4paper,margin=25mm,marginparwidth=28mm}
\usepackage{marginnote}
\usepackage{xcolor}
\usepackage[normalem]{ulem}
\setlength{\emergencystretch}{3em}
\usepackage{hyperref}
\hypersetup{colorlinks=true,linkcolor=black,urlcolor=black,pdfborder={0 0 0}}

% --- Batch metadata ---------------------------------------------------------
% Taken as they stand by the rendering script, which shows them at the head of
% the reading view. They fix what the file claims to cover. The macro names are
% the parent project's, so that the scripts stay shared; \folder here means the
% volume's slug — `fr-9115` — and \pages counts Gallica views.

\newcommand{\folder}[1]{\gdef\grothFolder{#1}}
\newcommand{\batch}[1]{\gdef\grothBatch{#1}}
\newcommand{\pages}[2]{\gdef\grothFirst{#1}\gdef\grothLast{#2}}
\newcommand{\foldertitle}[1]{\gdef\grothFoldertitle{#1}}

% The holder's own shelfmark, printed as they write it: « Français 9115 ».
\newcommand{\shelfmark}[1]{\gdef\grothShelfmark{#1}}

% The dating, copied verbatim from the holder's catalogue — never inferred
% here. « XIXe siècle » is the BnF's; a pass that dates a leaf from its content
% says so in a \note{}, as its own inference.
\newcommand{\dating}[1]{\gdef\grothDating{#1}}

% The status notice, printed in the title block and repeated as a diagonal
% watermark on every page of the PDF. The manuscripts are in the public
% domain; what this line declares is not a right but a quality — an
% unverified first pass by a machine. The skills write the line; a file
% without it gets no watermark, so leaving it out is visible in the output.
\newcommand{\watermark}[1]{\gdef\grothWatermark{#1}}

\gdef\grothFolder{?}\gdef\grothBatch{}\gdef\grothFirst{?}\gdef\grothLast{?}%
\gdef\grothFoldertitle{}\gdef\grothShelfmark{}\gdef\grothDating{}\gdef\grothWatermark{}

% --- The résumé -------------------------------------------------------------
\newenvironment{resume}
  {\small\setlength{\parskip}{0.45em}}
  {\par\medskip\hrule\medskip}

% --- Keywords ---------------------------------------------------------------
% In English, closing the modernised reading's résumé: the modern vocabulary
% under which to search for what the leaves construct. The manifest extracts
% them to tag the volume — this line is the single source of the tags.
\newcommand{\keywords}[1]{\par\smallskip\noindent
  {\footnotesize\textsc{Keywords} --- \textit{#1}}\par}

% --- The critical apparatus -------------------------------------------------

% The Gallica view number, in the margin. This is the anchor that lets a
% disagreement over a formula be settled by looking at *one* image rather than
% a volume of 750. The site uses it to turn the facsimile as the transcription
% is scrolled. A view is one scanned image: usually one side of a leaf.
\newcommand{\page}[1]{%
  \marginnote{\footnotesize\color{gray!70}\textsf{v.\,#1}}%
  \label{page:#1}\ignorespaces}

% The BnF's pencilled foliation, where the leaf carries it and a pass can read
% it: « 198r ». This is what the literature cites — « ff. 198r–208v » — and
% the one line that turns a citation into a view. Recorded, never inferred:
% a folio a pass cannot read is not written.
\newcommand{\folio}[1]{%
  \marginnote{\footnotesize\color{gray!50}\textsf{f.\,#1}}\ignorespaces}

% The modernised reading's coarser counterpart to \page{}: which views a
% section is drawn from.
\newcommand{\pagerange}[2]{%
  \marginnote{\footnotesize\color{gray!70}\textsf{v.\,#1--#2}}%
  \ignorespaces}

% Illegible: never guessed.
\newcommand{\ill}{\textcolor{red!60!black}{[\,\ldots\,]}}

% A doubtful reading: the word is given, and flagged as uncertain.
\newcommand{\uncertain}[1]{\uline{#1}}

% An editorial addition — a missing letter, an implied word.
\newcommand{\add}[1]{\textcolor{blue!50!black}{[#1]}}

% Struck out by the writer. What was crossed out is part of the document.
\newcommand{\struck}[1]{\sout{#1}}

% The transcriber's note, on the state of the page or the mathematics.
\newcommand{\note}[1]{\par\smallskip{\small\color{gray!60!black}\textit{[#1]}}\par\smallskip}

% A marginal note of hers, distinct from the body of the text.
\newcommand{\marginal}[1]{\par\smallskip\noindent\hspace{1em}%
  {\small\color{gray!60!black}#1}\par\smallskip}

% --- Title ------------------------------------------------------------------
% The title block is French, like the documents themselves.

\AddToHook{shipout/background}{%
  \ifx\grothWatermark\empty\else
    \begin{tikzpicture}[remember picture, overlay]
      \node[rotate=52, scale=3.4, align=center, text=gray!18,
            font=\sffamily\bfseries]
        at (current page.center) {\grothWatermark};
    \end{tikzpicture}%
  \fi}

\AtBeginDocument{%
  \begin{center}
    {\large\bfseries Manuscrits de mathématiciens (Gallica) ---
      \ifx\grothShelfmark\empty\grothFolder\else\grothShelfmark\fi}\\[2pt]
    {\itshape\grothFoldertitle}\\[4pt]
    {\small vues \grothFirst--\grothLast
      \ifx\grothBatch\empty\else\ (lot \grothBatch)\fi}\\[2pt]
    \ifx\grothDating\empty\else
      {\small\color{gray!60!black}Datation du catalogue : \grothDating}\\[2pt]
    \fi
    \ifx\grothWatermark\empty\else
      {\small\bfseries\color{red!45!gray}\grothWatermark}\\[2pt]
    \fi
    {\footnotesize\color{gray!60!black}Transcription établie d'après les images IIIF de Gallica.
     Source gallica.bnf.fr / Bibliothèque nationale de France.\\
     Les lectures incertaines sont soulignées, les illisibles notées [\,\ldots\,],
     les ajouts entre crochets ; \textsf{v.} numérote les vues de Gallica,
     \textsf{f.} la foliotation au crayon de la BnF.}
  \end{center}
  \vspace{1em}}

\endinput


\folder{fr-9115}
\shelfmark{Français 9115}
\batch{21}
\pages{401}{420}
\dating{1801-1900}
\watermark{Lecture automatique\\non vérifiée}
\foldertitle{Papiers de Sophie GERMAIN. — Recueil de dissertations et problèmes mathématiques et physiques.}

\begin{document}

\note{Numérotation des feuillets. Les feuillets de ce lot sont écrits
des deux côtés. En haut à droite des rectos, à l'encre, d'une plume plus
large et d'une écriture plus ronde, la série courante : 197 (vue 402),
198 (404), 199 (406), 200 (408), 201 (410), 202 (412), 203 (414), 204
(416), 205 (418), 206 (420), sans lacune, à la suite du 196 de la vue
400 ; les versos n'en portent aucun. Cette série occupe la place de la
foliotation de la Bibliothèque, mais elle est tracée à la plume, et
aucune foliotation au crayon n'a été lue ; aucun « f. » n'est donc
porté. En tête des pages, entre parenthèses, sa propre pagination : (4)
à (10) aux vues 401 à 407, puis (9) et (10) de nouveau aux vues 408 et
409, (11) barré (412), (12) (413), (13) barré (414), (14) (415), (15) à
(17) aux vues 418 à 420 ; les vues 410, 411 et 416 n'en portent pas.
Tous ces nombres ont été lus sur la vue ; aucun n'est déduit. La vue 417,
verso blanc du petit feuillet de la vue 416, n'est pas transcrite.}

\page{401}

on a demême pour le second plan que j'ai supposé passer
par le premier, le troisième et le quatrième point:
\begin{gather*}
x'=x,\quad y'=y,\quad z'=z;\qquad x''=x+dx,\quad y''=y,\quad z''=z+dz;\\
x'''=x+dx,\quad y'''=y+dy,\quad z'''=z+2dz+d
\end{gather*}
\[ A' = -z\,dy + (z+dz)\,dy + (z+2dz+d^2z).0 = dz\,dy \]
\[ B' = -x(dz+d^2z) + (x+dx)(2dz+d^2z) - (x+dx)\,dz = dx(dz+d^2z) \]
\[ C' = y.0 + y.dx - (y+dy)\,dx = -\,dy\,dx \]
\note{En tête, au milieu, « (4) » ; aucun nombre en haut à droite de ce
côté du feuillet. La première équation court jusqu'au bord du feuillet et
s'arrête sur « $+d$ » : la fin, que le calcul de $A'$ donne pour
$d^2z$, est perdue au bord ou n'a pas été écrite.}

Par conséquent, en se bornant au cas de \struck{ple} la surface plane, et
en supposant, comme l'ont fait tous les auteurs qui ont
traité des mouvemens de vibrations, que les differens points de
cette surface s'éloignent peu de leur situation \add{naturelle} pendant le mouvement,
lequel\struck{le} s'exécute tout entier dans la direction des $z$, il est
permis de négliger $dz$ vis avis de $dy$ et de $dx$. l'element
$ds$ de la surface se réduit à $dx\,dy$, et on trouve
\begin{align*}
A^2+B^2+C^2 &= dy^2(dz+d^2z)^2 + dz^2dx^2 + dx^2dy^2\\
&= dx^2dy^2 = ds^2\\
A'^2+B'^2+C'^2 &= dz^2dy^2 + dx^2(dz+d^2z)^2 + dy^2dx^2\\
&= dx^2dy^2 = ds^2
\end{align*}
\begin{align*}
&(AB'-A'B)^2+(AC'-A'C)^2+(BC'-B'C)^2\\
&= \{-dx\,dy(dz+d^2z)^2 + dz^2dx\,dy\}^2\\
&\quad + \{dy^2.dx(dz+d^2z) - dy^2dx\,dz\}^2\\
&\quad + \{dx^2dy\,dz - dx^2dy.(d^2z+dz)\}^2\\
&= dx^2dy^2\{((2dz+d^2z)\,d^2z\}^2 + \{dy.d^2z\}^2 + \{dx\,d^2z\}\\
&= ds^2(dx^2+dy^2)(d^2z)^2
\end{align*}
\begin{align*}
\text{sin}^2 e &= \frac{(AB'-A'B)^2+(AC'-A'C)^2+(BC'-B'C)^2}{(A^2+B^2+C^2)(A'^2+B'^2+C'^2)}\\
&= \frac{ds^2(dx^2+dy^2)(d^2z)^2}{ds^4} = \frac{(dx^2+dy^2)(d^2z)^2}{ds^2}
\end{align*}
\note{Les carrés sont écrits, comme partout sur ces feuillets, par un 2
placé au-dessus de la lettre ($dx$, $d^2z$, $e$) ; ils sont rendus par
des exposants ordinaires. Dans la troisième accolade, « $dx^2dy.(d^2z+dz)$ »
est écrit ainsi. À l'avant-dernière ligne, une parenthèse ouvrante est
doublée après l'accolade, et la dernière accolade n'a pas d'exposant : le
bord de la ligne touche la marge. Dans la dernière fraction, l'exposant de $d$ dans
« $(d^2z)^2$ » a la forme d'un 3 ; il est lu 2, comme partout ailleurs.}
\page{402}

\[ \text{sin}\,e = \frac{\sqrt{(dx^2+dy^2)}\,d^2z}{ds}; \]
et parceque l'angle $e$ est fort petit
\[ e = \frac{\sqrt{(dx^2+dy^2)}\,d^2z}{ds}, \qquad \delta e = \frac{\sqrt{(dx^2+dy^2)}\,\delta d^2z}{ds} \]
on voit aisément que la petite ligne $\beta$ doit être égale à la
distance entre le premier et le quatrième point, desorte que l'on a
\[ (z'''-z')^2+(y'''-y')^2+(x'''-x')^2 = (2dz+d^2z)^2+dy^2+dx^2 = dy^2+dx^2 = \beta^2 \]
ainsi $\int\!\!\int E\beta\,\delta e = \int\!\!\int E\beta^2\,\delta\,\frac{d^2z}{ds}$; et pour connoître la valeur
de la somme des momens, il ne s'agit plus que de déterminer
celle de la force $E$.
\note{En tête, au milieu, « (5) » ; en haut à droite « 197 », à l'encre,
de la grande écriture ronde de la série courante. Le signe de variation
est la lettre à haste bouclée que les passes précédentes ont rendue par
$\delta$ ; la double somme est écrite par deux s longs, rendue ici par
$\int\!\!\int$. La lettre de « la petite ligne $\beta$ » est un $\beta$ bas,
sans longue queue.}

Cette force peut être \struck{supposée} proportionnelle à la somme
des raisons inverses de l'un et l'autre rayons de courbure.
de plus, elle doit contenir un facteur constant dépendant de
la nature dela matière dela surface et de l'épaisseur
dela même surface. En adoptant cette hypothèse que
je crois être la plus convenable, on pourra donc \struck{prendre} faire
$E = b\tau^3\{\frac{1}{(r)}+\frac{1}{(r')}\}$ équation dans laquelle $b$ représente le
facteur dépendant de la rigidité naturelle, \uncertain{$\tau$} l'épaisseur,
et $(r)$, $(r')$ les deux rayons de courbures dela surface.
\note{La lettre de l'épaisseur, deux fois écrite, est un jambage
coiffé d'un trait tourné vers la droite, plus haut que le $r$ des rayons,
à la façon d'un $\Gamma$ ou d'un $\tau$ ; elle est rendue par $\tau$ avec doute. $b$
est souligné d'un trait. Dans « $(r')$ », le second
symbole est écrit sur un autre signe, repassé.}

J'ai déja dit que, pendant le mouvement, la surface vibrante
s'écartoit très peu \struck{du plan des $x$ et $y$} de sa position naturelle.
ainsi, puisqu'il ne s'agit ici que des surfaces planes, on pourra
\struck{prendre} \add{\struck{toujours} supposer} cette position \struck{pour} \add{\struck{très} voisine du} \struck{le} plan des $x$ et $y$. Il s'en suivra
que $\frac{dz}{dx}$ et $\frac{dz}{dy}$ seront des quantités très petites dont on pourra
\add{toujours} négliger les quarrés et les produits, et que l'on trouvera
\[ \frac{1}{(r)}+\frac{1}{(r')} = -\left(\frac{ddz}{dx^2}+\frac{ddz}{dy^2}\right) \]
\note{Le groupe biffé après « très peu » se lit « du plan des $x$ et
$y$ », les lettres $x$ et $y$ soulignées ; la lecture de « plan » est
malaisée sous la rature. Dans l'interligne au-dessus de « prendre cette
position pour le plan », les ajouts sont « toujours » biffé et
« supposer », puis « très » biffé et « voisine du ». Le « toujours »
ajouté au-dessus de « négliger » est traversé d'un trait fin qui peut
être une rature ; il est donné comme ajout. Le bas du feuillet, sous la
formule, est blanc.}
\page{403}

en intégrant par parties on a
\begin{align*}
\int\!\!\int E\beta^2\,\delta\frac{d^2z}{ds} &= \int E\beta^2\,\frac{d\,\delta z}{ds} - \int\!\!\int d.(E\beta^2).\frac{d\,\delta z}{ds}\\
&= \int E\beta^2\,\frac{d\,\delta z}{ds} - \int d.\frac{(E\beta^2)}{ds}.\delta z + \int\!\!\int dd.\frac{(E\beta^2)}{ds}.\delta z
\end{align*}
et comme les termes qui ne sont affectés que d'un seul signe
d'intégration, appartiennent aux limites, la somme des momens
des forces d'élasticité qui agissent dans toute l'étendue de
la surface, se reduit en fin à
\begin{align*}
\int\!\!\int dd.\frac{(E\beta^2)}{ds}.\delta z &= \int\!\!\int dd.\left\{\frac{-\left(\frac{ddz}{dx^2}+\frac{ddz}{dy^2}\right)(dx^2+dy^2)}{dx\,dy}\right\}b\tau^3.\delta z\\
&= -\int\!\!\int b\tau^3\left\{\frac{d^4z}{dx^4}+2\frac{d^4z}{dx^2dy^2}+\frac{d^4z}{dy^4}\right\}dx\,dy.\delta z
\end{align*}
\note{En tête, au milieu, « (6) » ; aucun nombre en haut à droite de ce
côté du feuillet, qui est le verso du feuillet 197 (vue 402). Dans la
dernière ligne, un 4 est écrit sous l'exposant 3 de $\tau$, contre
l'accolade ; ce qu'il marque n'est pas clair, et il n'est pas repris
dans la formule. La lettre $\tau$ est celle de l'épaisseur de la vue
402.}

En se bornant à la recherche de l'équation du mouvement
dela surface soumise à l'action de la seule force d'élasticité,
fesant abstraction du poids de cette surface, prenant la gravité
pour unité des forces accélératrices et se souvenant que le
mouvement se fait tout entier suivant les $z$, on aura
\begin{align*}
-\int\!\!\int b\tau^3\left\{\frac{d^4z}{dx^4}+2\frac{d^4z}{dx^2dy^2}+\frac{d^4z}{dy^4}\right\}dx\,dy.\delta z &= \int\!\!\int\frac{ddz}{dt^2}\,dm.\delta z\\
&= \int\!\!\int\frac{ddz}{dt^2}\,\tau.dx\,dy\,\delta z
\end{align*}
équation qui differenciée deux fois et divisée par $\tau\,dx\,dy\,\delta z$
donnera
\[ b\tau^2\left\{\frac{d^4z}{dx^4}+2\frac{d^4z}{dx^2dy^2}+\frac{d^4z}{dy^4}\right\}+\frac{ddz}{dt^2} = 0\ \ldots\ (A) \]
pour l'équation cherchée.

N.o 2. Lorsque le mouvement est régulier et comparable à
celui du pendule simple dont la longueur est $K$, on a,
comme l'on sait, en prenant pour unité des espaces parcourus
la hauteur \struck{delaquelle} \add{dont} un corps abandonné a lui même
tombe dans une seconde, $\left(\frac{ddz}{dt^2}\right) = -\frac{z}{K}$.
\note{Le $K$ du dénominateur est repassé d'une encre épaisse. Le bas du
feuillet, sous cette ligne, est blanc.}
\page{404}

ainsi l'équation de la surface élastique vibrante se réduit alors
à
\[ Kb\tau^2\left\{\frac{d^4z}{dx^4}+2\frac{d^4z}{dx^2dy^2}+\frac{d^4z}{dy^4}\right\} = z\ \ldots\ (B) \]
Elle peut s'abaisser au second degré; et, en fesant $Kb\tau^2 = h^4$
elle fournit les deux équations
\[ h^2\left\{\frac{d^2z}{dx^2}+\frac{d^2z}{dy^2}\right\} = z, \qquad h^2\left\{\frac{d^2z}{dy^2}+\frac{d^2z}{dx^2}\right\} = -z \]
\note{En tête, au milieu, « (7) » ; en haut à droite « 198 », à l'encre,
de la série courante. La lettre posée égale à $Kb\tau^2$ est un $h$ ;
dans les deux équations qui suivent, elle est tracée comme un b barré à
la haste, et c'est la même lettre. Le signe moins du second membre de la
seconde équation est serré contre $z$, peut-être ajouté.}

dont la dernière est, comme l'on voit, l'équation des mouvemens
réguliers des surfaces tendues. \struck{on peut} cette équation n'a point
d'intégrale générale en termes finis; ainsi tout cequ'on peut
espérer, est de trouver les intégrales particulières de l'une et
l'autre équations du second degré. \struck{Et} En joignant les deux
intégrales correspondante\struck{s}, on auroit, pour l'équation dela
surface\struck{s} elastique\struck{s}, une intégrale particulière delamême étendue
que celle \struck{que \ill{}} \add{dont jouit} l'integrale\struck{s} composante\struck{s} qui\struck{x} appartient à l'équation
dela surface\struck{s} tendue\struck{s} \struck{comporte pour ce genre de surface}.
\note{Elle supprime les pluriels de cette phrase en barrant d'une petite
croix chaque s final (« correspondantes », « surfaces », « elastiques »,
« integrales », « composantes », « quix » pour « qui », « tendues ») ;
ces croix sont données comme ratures. Au-dessus de « que » barré,
suivi d'une lettre barrée d'une croix qui n'a pas été lue, elle écrit
« dont jouit » ; la fin de la phrase,
« comporte pour ce genre de surface », est barrée d'un trait continu, et
la lecture de « ce genre » sous la rature est malaisée. « appartient »
est écrit au singulier sur la ligne.}

Il est clair que l'intégrale del'équation dela surface tendue
employée seule, peut egalement convenir à l'équation dela surface
élastique. On peut conclure delà que, \struck{quel que} \add{quelle que} soit la figure
que puisse prendre une surface tendue renfermée dans
des limites donnés, la même figure peut appartenir également
à une surface elastique renfermé dans les mêmes limites.
\note{Le bas du feuillet, sous cette ligne, est blanc.}
\page{405}

\section{§. 2. Intégration del'équation differentielle dela surface elastique vibrante.}

N.o 3. La condition \struck{N.o 2} $\left(\frac{ddz}{dt^2}\right) = -\frac{z}{K}$ donnée §.1 N.o 2. fournit
l'équation $z = \Pi\,\text{sin.}\{t.\sqrt{\tfrac{1}{K}}+\zeta\}$ dans laquelle \struck{$\Pi$ et} $\zeta$
représente\struck{nt} \struck{des} \add{une} constante\struck{s} et $\Pi$ une fonction de $x$ $y$ et constantes.
C'est la nature de cette fonction que détermine\struck{nt} les
differentes intégrales particulières de l'équation \struck{dela surface
elastique} (B). Il est clair que cette équation qui n'est qu'une
\struck{transformation} \add{suite} de l'équation (A), doit fournir pour $z$ des
valeurs qui contiendront toutes le facteur $\text{sin}(t.\sqrt{\tfrac{1}{K}}+\zeta)$
relatif au tems. Ainsi je me dispenserai, dans la suite,
d'écrire ce facteur qui devra toujours etre sousentendu
dans chaque valeur de $z$.
\note{En tête, au milieu, « (8) » ; aucun nombre en haut à droite de ce
côté du feuillet, verso du feuillet 198 (vue 404). Le titre du § 2 est
souligné au milieu d'un trait court. Au-dessus de « condition », un
petit groupe barré se lit peut-être « N.o 2 ». La lettre de la fonction
est un grand $\Pi$ aux jambages hauts, rendu par $\Pi$. Les fins de mots
barrées de « représentent », « des constantes », « déterminent » sont
barrées d'une croix ou d'un trait, comme à la vue 404.}

Pour trouver les intégrales de l'équation (B), \add{il est permis} \struck{on peut}
davoir recours à differentes conditions que l'on peut imaginer
entre les variables $z$, $x$ et $y$.

Je supposerai d'abord que \struck{que} l'équation $\frac{ddz}{dx\,dy} = 0$ soit
\struck{remplie} \add{satisfaite}, c'est adire que $z$ se compose de la somme
de deux fonctions séparées des variables $x$ et $y$. on pourra
prendre alors
\[ z = \ldots\left\{\begin{array}{l}
\mathcal{A}\left\{a\,e^{\frac{\omega x}{A}}+b\,e^{-\frac{\omega x}{A}}+c\,\text{sin.}\frac{\omega x}{A}+d\,\text{cos.}\frac{\omega x}{A}\right\}\\[1ex]
+B\left\{a\,e^{\frac{\omega y}{A}}+b\,e^{-\frac{\omega y}{A}}+c\,\text{sin.}\frac{\omega y}{A}+d\,\text{cos.}\frac{\omega y}{A}\right\}
\end{array}\right\}\ \ldots\ (C) \]
les sons simples qui se rapportent aux \struck{cette} manières
de vibrer renfermées dans cette intégrale, seront
exprimés par la formule $\sqrt{\tfrac{1}{K}}.\frac{1}{\pi} = \frac{\omega\omega\tau\sqrt{b}}{\pi.AA}$ $\pi$ représentant,
\note{Le mot barré sous « satisfaite » se lit « remplie » avec doute. « $z
= \ldots$ » : les points sont sur la page, pour le facteur du temps
sousentendu. Le coefficient de la première ligne est une majuscule
cursive bouclée, distincte du $A$ des dénominateurs et du $B$ de la
seconde ligne ; elle est rendue par $\mathcal{A}$, la lecture restant
douteuse (la lettre a aussi la forme d'un B cursif). Dans la formule des sons, « $\sqrt{\frac{1}{K}}.\frac{1}{\pi}$ »
est écrit au-dessus du signe $=$, en premier membre. La phrase se
poursuit à la vue 406.}
\page{406}

comme à l'ordinaire, la demie circonference.

Dans l'équation précédente, je n'ai écrit que le premier
terme dela valeur de $z$, et j'en userai demême pour \add{toutes} les valeurs
de $z$ que je donnerai par \uncertain{lasuite} \add{car} le but de mes recherches
étant la détermination des differentes figures nodales et des
sons qui leur correspondent, je n'ai pas à traiter de
la coexistence des sons qui empêcheroit la formation
de ces figures et qui ne présente d'ailleurs pour les
surfaces élastiques aucune difficulté de plus que pour les
differens corps sonores dont les géomètres se sont occupés
jusqu'à présent.
\note{En tête, au milieu, « (9) » ; en haut à droite « 199 », à l'encre,
de la série courante. La première ligne achève la phrase de la vue 405
(« $\pi$ représentant, comme à l'ordinaire, la demie circonference »).
« lasuite » est chargé et en partie repassé, sous l'ajout « car » écrit
au-dessus ; le mot n'est pas nettement barré.}

N.o 4. On peut encore supposer qu'il existe une periodicité
dans les valeurs des coordonnées $x$ et $y$, c'est-adire que ces
variables augmentées l'une et l'autre d'une quantité
constante, mais qui peut etre differente pour chacune d'elles,
redonnent les mêmes valeurs de $z$. Cette condition peut
etre exprimée par l'équation $\left(\frac{ddz}{dx^2}\right) = \left(\frac{ddz}{dy^2}\right)\frac{n^2}{m^2}$ dela
forme de celle aux cordes vibrantes. Il en résultera
\[ z = \ldots\ \mathcal{A}\,\text{sin.}\frac{(\pi.nx)}{A}\ \text{sin}\frac{(\pi.my)}{A}\ \ldots\ldots\ (D) \]
et les sons simples qui \struck{sont} \add{se rapportent aux manières de vibrer} renfermées dans cette integrale
seront exprimés par la formule $\sqrt{\tfrac{1}{K}}.\frac{1}{\pi} = \frac{\pi.(m^2+n^2)\tau\sqrt{b}}{A^2}$.

La supposition de $\mathcal{A}$ constant n'est pas la seule
qui puisse satisfaire à la condition $\left(\frac{ddz}{dx^2}\right) = \left(\frac{ddz}{dy^2}\right)\frac{n^2}{m^2}$
On peut encore faire
\note{« cordes » est écrit sur un autre mot, d'une encre plus épaisse. Le
second dénominateur de (D) est un $A$ traversé d'un trait oblique, qui
peut être un accent ($A'$) ou une rature. Le coefficient $\mathcal{A}$ de
(D) et de la phrase qui suit est la même majuscule cursive bouclée qu'en
tête de (C) à la vue 405. Dans la condition, l'exposant de $n$ est
écrit au-dessus de la lettre. La page s'arrête sur « On peut encore
faire », le bas du feuillet blanc ; la suite est à la vue 407.}
\page{407}

\begin{align*}
z = \ldots\ \text{sin}\,\pi.\frac{nx}{A}.\ \text{sin.}\,\pi.\frac{my}{A}\Bigl(\,&a'\Bigl\{\text{cos.}\,\pi.\frac{nx}{A} + \frac{\Pi+3}{2.3}\,\text{cos.}^3\pi.\frac{nx}{A}\\
&\qquad + \frac{(\Pi+3)(\Pi+3.5)}{2.3.4.5}\,\text{cos}^5\pi\frac{nx}{A} + \&c\Bigr\}\\
&+b'\Bigl\{\frac{2}{\Pi} + \text{cos}^2\pi.\frac{nx}{A} + \frac{\Pi+8}{3.4}\,\text{cos}^4\pi.\frac{nx}{A}\\
&\qquad + \frac{(\Pi+8)(\Pi+3.8)}{3.4.5.6}\,\text{cos}^6\pi.\frac{nx}{A} + \&c\Bigr\}\\
\times\ &A'\Bigl\{\text{cos}\,\pi.\frac{my}{A} + \frac{\Pi+3}{2.3}\,\text{cos}^3\pi.\frac{my}{A}\\
&\qquad + \frac{(\Pi+3)(\Pi+3.5)}{2.3.4.5}\,\text{cos}^5\pi.\frac{my}{A} + \&c\Bigr\}\\
&+B'\Bigl\{\frac{2}{\Pi} + \text{cos}^2\pi.\frac{my}{A} + \frac{\Pi+8}{3.4}\,\text{cos}^4\pi.\frac{my}{A}\\
&\qquad + \frac{(\Pi+8)(\Pi+3.8)}{3.4.5.6}\,\text{cos}^6\pi\frac{nx}{A}\Bigr\}\Bigr)\ \ldots\ (E)
\end{align*}
\note{En tête, au milieu, « (10) » ; aucun nombre en haut à droite de ce
côté du feuillet, verso du feuillet 199 (vue 406). La formule (E) est
écrite sur toute la largeur, en quatre lignes réunies par de grandes
accolades, « (E) » et une ligne de points à gauche ; sa disposition est
simplifiée ici. Le même signe, un $\pi$ aux jambages hauts, sert pour la
demi-circonférence dans les arguments des sinus et cosinus et pour une
quantité qui entre dans les coefficients (« $\Pi+3$ », « $\frac{2}{\Pi}$ »,
et « $\Pi-1$ » dans la formule des sons) ; il est rendu par $\pi$ dans le
premier emploi et par $\Pi$ dans le second, la distinction étant du
transcripteur et non de la page. Le point entre deux chiffres est son
signe de produit : « $3.5$ » et « $3.8$ » valent 15 et 24. Le dernier
terme de la ligne en $B'$ porte « $\frac{nx}{A}$ » là où les autres
termes de la ligne portent « $\frac{my}{A}$ », et n'a pas de « \&c » ;
c'est la page. Dans les lignes en $y$, le dénominateur $A$ est plusieurs
fois traversé d'un trait (comme à la vue 406).}

et les sons simples qui correspondent aux manières de vibrer \struck{qui
sont} renfermées dans cette intégrale, seront exprimés par la
formule $\sqrt{\tfrac{1}{K}}.\frac{1}{\pi} = \frac{\pi.(\Pi-1)(m^2+n^2)\tau\sqrt{b}}{A^2}$.

Mais si on fait
\begin{align*}
z = \ldots\ \text{sin.}\,\pi.\frac{mx}{A}\ \text{sin}\,\pi.\frac{my}{A}\Bigl(\,&a'\Bigl\{\text{cos}\,\pi\frac{mx}{A} + \frac{(\Pi'+3)}{2.3}\,\text{cos}^3\pi.\frac{mx}{A}\\
&\qquad + \frac{((\Pi'+3)((\Pi'+3.5)}{2.3.4.5}\,\text{cos}^5\pi.\frac{mx}{A} + \&c\Bigr\}\\
&+b'\Bigl\{\frac{2}{\Pi'} + \text{cos}^2\pi.\frac{mx}{A} + \frac{(\Pi'+8)}{3.4}\,\text{cos}^4\pi.\frac{mx}{A}\\
&\qquad + \frac{((\Pi'+8)((\Pi'+3.8)}{3.4.5.6}\,\text{cos}^6\pi\frac{mx}{A}\Bigr\}\\
\times\ &A'\Bigl\{\text{cos.}\,\pi.\frac{my}{A} + \frac{\Pi'+3}{2.3.}\,\text{cos}^3\pi.\frac{my}{A}\\
&\qquad + \frac{(\Pi'+3)(\Pi'+3.5)}{2.3.4.5}\,\text{cos}^5\pi.\frac{mx}{A} + \&c\Bigr\}\\
&+B'\Bigl\{\frac{2}{\Pi'} + \text{cos}^2\pi.\frac{mx}{A} + \frac{(\Pi'+8)}{3.4}\,\text{cos.}^4\pi.\frac{mx}{A}\\
&\qquad + \frac{(\Pi'+8)(\Pi'+3.8)}{3.4.5.6}\,\text{cos}^6\pi\frac{mx}{A}\ \struck{+\&c}\Bigr\}\Bigr)\ \ldots\ (F)
\end{align*}
\note{Dans les deux premières lignes de (F), chaque $\Pi$ des
coefficients est suivi d'un petit groupe barré de croix (« $-$ » et deux
lettres, qui n'ont pas été lues) au-dessus duquel est posé l'accent ;
dans les deux dernières lignes, l'accent est simplement ajouté au-dessus
de $\Pi$. La correction fait donc partout lire $\Pi'$, qui est donné ici.
Les parenthèses doublées « $((\Pi'+3)$ » sont celles de la page. Au
dénominateur de $\frac{2}{\Pi'}$ de la deuxième ligne, et sous
« $\text{cos}^2\pi.\frac{mx}{A}$ », le trait et le $A$ sont repassés
d'une encre épaisse. Dans les lignes en $y$, plusieurs arguments portent
$mx$ où l'on attend $my$ ; c'est la page. « $+\&c$ » est barré à la fin
de la dernière ligne.}

Les sons simples seront exprimés par la formule $\sqrt{\tfrac{1}{K}}.\frac{1}{\pi} = \frac{\pi.m^2(\Pi+\Pi'-2)\tau\sqrt{b}}{A^2}$
et au lieu dela condition $\left(\frac{ddz}{dx^2}\right) = \left(\frac{ddz}{dy^2}\right)\frac{n^2}{m^2}$ on aura celle ci:
\[ \left(\frac{ddz}{dx^2}\right) = \left(\frac{ddz}{dy^2}\right)\frac{\Pi'\,\ill\,-1}{\Pi'-1} \]
\note{Dans la formule des sons, une lettre noircie d'encre précède
« $\Pi+\Pi'-2$ » dans la parenthèse. Au numérateur de la dernière
condition, après $\Pi'$, un groupe est barré (un signe et une lettre
surmontée d'un petit x) ; il n'a pas été lu, et ce que la correction
laisse au numérateur n'est pas sûr.}

N.o 5. En prenant $z = z'z''$, la surface sera encore periodique
dans le sens des $x$ et dans celui des $y$ lorsque l'on aura
la condition $\left(\frac{ddz'}{dx^2}\right) = \left(\frac{ddz'}{dy^2}\right)\frac{n^2}{m^2}$; et alors les differentes
valeurs de $z''$ détermineront la nature des periodes.
\note{L'accent du premier facteur de « $z'z''$ » est petit et serré
contre la lettre. Le bas du feuillet est blanc.}
\page{408}

Mais si on fait
\begin{align*}
z = \ldots\ \text{sin.}\frac{(\pi.mx)}{A}\ \text{sin}\frac{(\pi.my)}{A}\Bigl(\,&a'\Bigl\{\text{cos.}\frac{(\pi.mx)}{A} + \frac{\Pi+3}{2.3}\,\text{cos.}^3\frac{(\pi.mx)}{A}\\
&\qquad + \frac{(\Pi+3)(\Pi+3.5)}{2.3.4.5}\,\text{cos}^5\frac{(\pi.mx)}{A} + \&c\Bigr\}\\
&+b'\Bigl\{\frac{2}{\Pi} + \text{cos}^2\frac{(\pi.mx)}{A} + \frac{\Pi+8}{3.4}\,\text{cos}^4\frac{(\pi.mx)}{A}\\
&\qquad + \frac{(\Pi+8)(\Pi+3.8)}{3.4.5.6}\,\text{cos.}^6\frac{(\pi.mx)}{A} + \&c\Bigr\}\\
\times\ &A'\Bigl\{\text{cos.}\frac{(\pi.my)}{A} + \frac{\Pi'+3}{2.3}\,\text{cos.}^3\frac{(\pi.my)}{A}\\
&\qquad + \frac{(\Pi'+3)(\Pi'+3.5)}{2.3.4.5}\,\text{cos.}^5\frac{(\pi.my)}{A} + \&c\Bigr\}\\
&+B'\Bigl\{\frac{2}{\Pi'} + \text{cos}^2\frac{(\pi.my)}{A} + \frac{\Pi'+8}{3.4}\,\text{cos.}^4\frac{(\pi.my)}{A}\\
&\qquad + \frac{(\Pi'+8)(\Pi'+3.8)}{3.4.5.6}\,\text{cos.}^6\frac{(\pi.my)}{A} + \&c\Bigr\}\Bigr)\ \ldots\ (F')
\end{align*}
\note{En tête, vers la droite, « (9) » ; en haut à droite « 200 », à
l'encre, de la série courante. Ce feuillet est une mise au net de la
seconde moitié de la page (10) de la vue 407 : la même formule, les
accents de $\Pi'$ écrits d'emblée dans les lignes en $y$, les arguments
mis entre parenthèses, et la formule étiquetée « (F') ». Le chiffre de
la tête se lit nettement 9, bien que la page (9) soit déjà à la vue
406 ; il est relevé tel quel. Sur la notation $\pi$ et $\Pi$, et sur le
point de produit, voir la note de la vue 407.}

Les sons simples seront exprimés par la formule $\sqrt{\tfrac{1}{K}}.\frac{1}{\pi} = \frac{\pi.m^2(\Pi+\Pi'-2)\tau\sqrt{b}}{A^2}$
et au lieu de la condition $\frac{ddz}{dx^2} = \frac{n^2}{m^2}\frac{ddz}{dy^2}$ on aura celle ci: $\frac{ddz}{dx^2} = \frac{\Pi-1}{\Pi'-1}.\frac{ddz}{dy^2}$.

N.o 5. En prenant $z = z'z''$, la surface sera encore périodique
dans le sens des $x$ et dans celui des $y$ lorsque l'on aura la condition
$\frac{ddz'}{dx^2} = \frac{n^2}{m^2}.\frac{ddz'}{dy^2}$; et alors les differentes valeurs de $z''$ détermineront la
nature des périodes. Ainsi lorsque l'on fera $z' = \text{sin}\frac{(\pi.xm)}{A}\ \text{sin}\frac{(\pi ym)}{A}$, on pourra
prendre
\begin{align*}
z = \ldots\ \text{sin}\frac{(\pi mx)}{A}\ \text{sin}\frac{(\pi my)}{A}\Bigl[\\ &a'\Bigl\{\text{cos}\frac{(\pi mx)}{A} + \frac{\Pi+3}{2.3}\,\text{cos}^3\frac{(\pi mx)}{A} + \ldots\ldots\\
&\quad \frac{(\Pi+3)(\Pi+3.5)\ldots(\Pi+4n^2-1)}{2.3.4.5\ldots 2n+1}\,\text{cos}^{2n+1}\frac{(\pi mx)}{A} + \&c\Bigr\}\\
&+b'\Bigl\{\frac{2}{\Pi} + \text{cos}^2\frac{(\pi mx)}{A} + \frac{\Pi+8}{3.4}\,\text{cos}^4\frac{(\pi mx)}{A} + \ldots\ldots\\
&\quad \frac{(\Pi+8)(\Pi+3.8)\ldots(\Pi+(2n-1)^2-1)}{3.4.5.6\ldots 2n}\,\text{cos}^{2n}\frac{(\pi mx)}{A} + \&c\Bigr\}\\
\times\ &A'\Bigl\{\text{cos}\frac{(\pi my)}{A} + \frac{\Pi'+3}{2.3}\,\text{cos}^3\frac{(\pi my)}{A} + \ldots\ldots\\
&\quad \frac{(\Pi'+3)(\Pi'+3.5)\ldots(\Pi'+4n^2-1)}{2.3.4.5\ldots 2n+1}\,\text{cos}^{2n+1}\frac{(\pi my)}{A} + \&c\Bigr\}\\
&+B'\Bigl\{\frac{2}{\Pi'} + \text{cos}^2\frac{(\pi my)}{A} + \frac{\Pi'+8}{3.4}\,\text{cos}^4\frac{(\pi my)}{A} + \ldots\ldots\\
&\quad \frac{(\Pi'+8)(\Pi'+3.8)\ldots(\Pi'+(2n-1)^2-1)}{3.4.5.6\ldots 2n}\,\text{cos}^{2n}\frac{(\pi my)}{A} + \&c\Bigr\}\\
&+a''\Bigl\{\text{cos}\frac{(\pi mx)}{A} + \frac{\Pi''+3}{2.3}\,\text{cos}^3\frac{(\pi mx)}{A} + \ldots\ldots\\
&\quad \frac{(\Pi''+3)(\Pi''+3.5)\ldots(\Pi''+4n^2-1)}{2.3.4.5\ldots 2n+1}\,\text{cos}^{2n+1}\frac{(\pi mx)}{A} + \&c\Bigr\}\\
&+b''\Bigl\{\frac{2}{\Pi''} + \text{cos}^2\frac{(\pi mx)}{A} + \frac{\Pi''+8}{3.4}\,\text{cos}^4\frac{(\pi mx)}{A} + \ldots\ldots\\
&\quad \frac{(\Pi''+8)(\Pi''+3.8)\ldots(\Pi''+(2n-1)^2-1)}{3.4.5.6\ldots 2n}\,\text{cos}^{2n}\frac{(\pi mx)}{A} + \&c\Bigr\}\\
&+A''\Bigl\{\text{cos}\frac{(\pi my)}{A} + \frac{\Pi''+3}{2.3}\,\text{cos}^3\frac{(\pi my)}{A} + \ldots\ldots\\
&\quad \frac{(\Pi''+3)(\Pi''+3.5)\ldots(\Pi''+4n^2-1)}{2.3.4.5\ldots 2n+1}\,\text{cos}^{2n+1}\frac{(\pi my)}{A} + \&c\Bigr\}\\
&+B''\Bigl\{\frac{2}{\Pi''} + \text{cos}^2\frac{(\pi my)}{A} + \frac{\Pi''+8}{3.4}\,\text{cos}^4\frac{(\pi my)}{A} + \ldots\ldots\\
&\quad \frac{(\Pi''+8)(\Pi''+3.8)\ldots(\Pi''+(2n-1)^2-1)}{3.4.5.6\ldots 2n}\,\text{cos}^{2n}\frac{(\pi my)}{A} + \&c\Bigr\}\\
&\Bigr]\ \ldots\ (G)
\end{align*}
pourvu que $\Pi$, $\Pi'$ et $\Pi''$ ayent entre elles la relation exprimée par l'équation
$\Pi+\Pi' = \Pi''$,

Les sons simples seront comme dans les cas compris dans l'intégrale (F),
\note{La formule (G) occupe huit lignes réunies par un grand crochet, avec
« (G) » et une ligne de points à gauche ; sa disposition est simplifiée
ici. Chaque ligne donne les deux ou trois premiers termes, une ligne de
points, puis le terme général. L'exposant de « $4n^2$ » et celui de
« $(2n-1)^2$ » sont écrits au-dessus de la parenthèse ou de la lettre. Le
second accent de $\Pi''$ dans « $\Pi+\Pi' = \Pi''$ » est repassé d'une
encre épaisse sur une autre lettre. La phrase finale s'arrête sur une
virgule au bas du feuillet ; la suite est à la vue 409.}
\page{409}

c'est-a-dire que l'on aura toujours $\sqrt{\tfrac{1}{K}}.\frac{1}{\pi} = \frac{\pi.m^2(\Pi+\Pi'-2)}{A^2}$.

Lorsque l'on prendra $z' = \text{sin}\frac{(\pi nx)}{A}\ \text{sin.}\frac{(\pi.my)}{A}$, on pourra faire
\begin{align*}
z = \ldots\ \text{sin}\frac{(\pi.nx)}{A}\ \text{sin}\frac{(\pi.my)}{A}\Bigl\{\,&a'\Bigl\{\text{cos}\frac{(\pi nx)}{A} + \frac{(2-m^2+3}{2.3}\,\text{cos}^3\frac{(\pi nx)}{A}\\
&\qquad + \frac{((2-m^2+3)((2-m^2+15)}{2.3.4.5}\,\text{cos}^5\frac{(\pi nx)}{A} + \&c\Bigr\}\\
&+b'\Bigl\{\frac{2}{2-m^2} + \text{cos}^2\frac{(\pi nx)}{A} + \frac{(2-m^2+8}{3.4}\,\text{cos}^4\frac{(\pi nx)}{A}\\
&\qquad + \frac{((2-m^2+8)((2-m^2+24)}{3.4.5.6}\,\text{cos}\ \ldots\Bigr\}\\
&+A'\Bigl\{\text{cos}\frac{(\pi my)}{A} + \frac{(2-n^2+3}{2.3}\,\text{cos}^3\frac{(\pi.my)}{A}\\
&\qquad + \frac{((2-n^2+3)((2-n^2+15)}{2.3.4.5}\,\text{cos}^5\frac{(\pi my)}{A}\ \ldots\Bigr\}\\
&+B'\Bigl\{\frac{2}{2-n^2} + \text{cos}^2\frac{(\pi my)}{A} + \frac{(2-n^2+8}{3.4}\,\text{cos}^4\frac{(\pi my)}{A}\\
&\qquad + \frac{((2-n^2+8)((2-n^2+24)}{3.4.5.6}\ \ldots\Bigr\}\ \ldots\ (H)
\end{align*}
\note{En tête, au milieu, « (10) » ; aucun nombre en haut à droite de ce
côté du feuillet, verso du feuillet 200 (vue 408). Le chiffre de la tête
répète celui de la vue 407 ; il est relevé tel quel. La première ligne
achève la phrase de la vue 408 (« Les sons simples seront comme dans les
cas compris dans l'intégrale (F), c'est-a-dire… »). Dans (H), les
coefficients s'écrivent « $(2-m^2+3$ » avec la parenthèse ouverte et non
fermée, l'exposant 2 posé au-dessus de $m$ ; le premier chiffre de ces
groupes est chargé d'encre dans plusieurs termes et se lit 2 sans
certitude. « $15$ » et « $24$ » sont écrits en chiffres, là où les
formules des vues 407 et 408 écrivent « $3.5$ » et « $3.8$ ». Les termes
de droite des lignes en $b'$, $A'$ et $B'$ sont coupés par le bord du
feuillet ou restent inachevés ; les points marquent ce qui manque. Un
long trait vertical, sans doute accidentel, traverse les quatre lignes
de la formule.}

Les sons simples qui appartiennent aux manières de vibrer renfermées dans
cette dernière intégrale, seront exprimés par la formule $\sqrt{\tfrac{1}{K}}.\frac{1}{\pi} = \frac{\pi.m^2n^2\tau\sqrt{b}}{A^2}$.

N.o 6. En prenant encore $z = z'z''$, et fesant $z' = \text{sin}(t.\sqrt{\tfrac{1}{K}}+\zeta)\,\text{sin}\frac{(\pi nx)}{A}$,
on aura toujours une équation de la forme de celle aux cordes
vibrantes savoir $\frac{ddz'}{dt^2} = \frac{K}{n^2}.\frac{ddz'}{dx^2}$. Elle indique que la surface sera
périodique dans le sens des $x$. On pourra faire
\begin{align*}
z = \ldots\ \text{sin}\frac{(\pi nx)}{A}\Bigl\{\,&a'\Bigl\{\text{cos.}\frac{(\pi.nx)}{A} + \frac{2-m^2+3}{2.3}\,\text{cos}^3\frac{(\pi nx)}{A} + \&c\Bigr\}\\
&+ b'\Bigl\{\frac{2}{2-m^2} + \text{cos}^2\frac{(\pi nx)}{A} + \frac{2-m^2+8}{3.4}\,\text{cos}^4\frac{(\pi nx)}{A} + \&c\Bigr\}\\
&+\ \mathfrak{A}\,\text{sin.}\frac{(\pi\,\uncertain{may})}{A}\Bigr\}\ \ldots\ldots\ (I)
\end{align*}
\note{Dans (I), un trait épais barre le signe qui précède « $2-m^2$ » dans
les deux coefficients (une parenthèse ouvrante, semble-t-il). Le
coefficient de la dernière ligne est une majuscule gothique ou très
ornée, rendue par $\mathfrak{A}$ ; l'argument de son sinus se lit
« $\pi\,may$ » avec doute, la lettre du milieu pouvant être une autre
lettre écrite sur $m$.}

mais cette surface ne jouira pas de la même propriété dans
le sens des $y$, puisque cette dernière variable n'entre pas dans l'équation
aux cordes vibrantes. Les sons simples seront exprimés par la formule $\sqrt{\tfrac{1}{K}}.\frac{1}{\pi} = \struck{\frac{\pi m^2n^2\tau\sqrt{b}}{A^2}}$
$= \pi.n^2(m^2+1)$,
\note{La première valeur, qui répète celle de (H), est barrée ; la
seconde, « $= \pi.n^2(m^2+1)$ », est écrite au-dessous, en plus petit,
sans le facteur $\frac{\tau\sqrt{b}}{A^2}$.}

Les differentes intégrales que je viens de donner, conviennent toutes, exc\add{epté}
la première, à l'équation differentielle des surfaces tendues; et c'est
\add{à} quoi elles doivent leur simplicité, car si on vouloit avoir les intégra\add{les}
qui n'appartiennent qu'à l'équation differentielle des surfaces élastiq\add{ues}
on trouveroit non seulement les termes contenans $e^{\pm\sqrt{-1}.\omega}$ mais encore
d'autres termes ou entreroient de la même manière, les exponentielles
$e^{\pm\omega}$; et alors l'angle $\omega$ ne pourroit être déterminé que par
des calculs d'aproximation; desorte que non seulement les intégrales
\note{Les fins de ligne « exc », « intégra », « élastiq » passent sous le
bord droit du feuillet, et « à » manque devant « quoi » au début de la
ligne suivante ; ce que le sens rend certain est ajouté entre crochets. L'exposant de
« $e^{\pm\sqrt{-1}.\omega}$ » est écrit petit au-dessus de la ligne. La
phrase se poursuit au feuillet suivant.}
\page{410}

\note{Feuillet de calculs, sans une ligne de prose, écrit en travers : la
vue le montre couché, et le « 201 » à l'encre de la série courante se lit
en le redressant, en haut à droite du feuillet ainsi tourné ; aucune
pagination de sa main. L'écriture est plus rapide que celle des pages
(5) à (10), avec des surcharges. Les formules sont données dans l'ordre
de la page, de haut en bas ; leur disposition sur plusieurs lignes, en
accolades et crochets, est simplifiée. Même notation qu'aux vues 407 à
409 ($\pi$ dans les arguments, $\Pi$, $\Pi'$ dans les coefficients).}

\begin{align*}
z = \text{sin}(\zeta + t\sqrt{\tfrac{1}{K}})\ \text{sin}.\pi.\frac{xm}{A}\ \text{sin}.\pi.\frac{yn}{A}\Bigl[\\ &a'\Bigl\{\text{cos}.\pi.\frac{mx}{A} + \frac{(\Pi-\Pi')+3}{2.3}\,\text{cos}^3\pi.\frac{mx}{A}\\
&\qquad + \frac{((\Pi-\Pi')+3)((\Pi-\Pi')+3.5)}{2.3.4.5}\,\text{cos}^5\pi.\frac{mx}{A} + \&c\Bigr\}\\
&+ b'\Bigl\{\frac{2}{\Pi-\Pi'} + \text{cos}^2\pi\frac{mx}{A} + \frac{(\Pi-\Pi')+8}{3.4}\,\text{cos}^4\pi\frac{mx}{A}\\
&\qquad + \frac{((\Pi-\Pi')+8)((\Pi-\Pi')+3.8)}{3.4.5.6}\,\text{cos}^6\pi\frac{mx}{A} + \&c\Bigr\}\\
\times\ &A'\Bigl\{\text{cos}.\pi.\frac{ny}{A} + \frac{\Pi'+3}{2.3}.\text{cos}^3\pi.\frac{ny}{A}\\
&\qquad + \frac{(\Pi'+3)(\Pi'+3.5)}{2.3.4.5}.\text{cos}^5\pi\frac{ny}{A} + \&c\Bigr\}\\
&+ B'\Bigl\{\frac{2}{\Pi'} + \text{cos}^2\frac{ny}{A} + \frac{\Pi'+8}{3.4}\,\text{cos}^4\frac{ny}{A}\\
&\qquad + \frac{(\Pi'+8)(\Pi'+3.8)}{3.4.5.6}\,\text{cos}^6\pi\frac{ny}{A} + \&c\Bigr\}\\
&+a\Bigl\{\text{cos}\,\pi\frac{\sqrt{m^2+n^2}\,x}{A} + \frac{\Pi+3}{2.3}\,\text{cos}^3\pi\frac{\sqrt{m^2+n^2}\,x}{A}\\
&\qquad + \frac{(\Pi+3)(\Pi+3.5)}{2.3.4.5}\,\text{cos}^5\pi.\frac{\sqrt{m^2+n^2}}{A}x + \&c\Bigr\}\\
&+ b\Bigl\{\text{cos}^2\pi\frac{\sqrt{m^2+n^2}\,x}{A} + \frac{\Pi+8}{3.4}\,\text{cos}^4\pi.\frac{\sqrt{m^2+n^2}}{A}x\\
&\qquad + \frac{(\Pi+8)(\Pi+3.8)}{3.4.5.6}\,\text{cos}^6\ldots\Bigr\}\\
&+A\Bigl\{\text{cos}\,\pi\sqrt{m^2+n^2}.\frac{y}{A} + \frac{\Pi+3}{2.3}.\text{cos}^3\pi\sqrt{m^2+n^2}.\frac{y}{A}\\
&\qquad + \frac{(\Pi+3)(\Pi+3.5)}{2.3.4.5}\,\text{cos}^5\pi\sqrt{m^2+n^2}\frac{y}{A} + \&c\Bigr\}\\
&+ B\Bigl\{\text{cos}^2\pi.\sqrt{m^2+n^2}.\frac{y}{A} + \frac{\Pi+8}{3.4}\,\text{cos}^4\pi\sqrt{m^2+n^2}\frac{y}{A}\\
&\qquad + \frac{(\Pi+3)(\Pi+3.4)}{3.4.5.6}\,\text{cos}^6\pi\sqrt{m^2+n^2}\frac{y}{A}\ldots\Bigr\}\\
&+\ \mathfrak{A}\Bigl(\frac{b+B}{\Pi}\Bigr)
\end{align*}
\note{Au bout de la ligne en $b'$, sous le dernier terme, « $+\&c$ » est
écrit à part. Dans la ligne en $b$, « $\frac{\Pi+8}{3.4}$ » a le 8 écrit
sur un autre chiffre. Dans le dernier terme de la ligne en $B$, la page
porte « $(\Pi+3)(\Pi+3.4)$ » là où les autres lignes en $b$ portent
« $(\Pi+8)(\Pi+3.8)$ » ; c'est la page. Les fins des lignes en $b$ et en
$B$ sont coupées au bord droit. Le coefficient de la dernière ligne est la
majuscule ornée de la vue 409, rendue par $\mathfrak{A}$.}

\[ \frac{dz}{dx} = -\pi\frac{m}{A}\,\text{cos}.\pi\frac{mx}{A}\{\ldots\ldots \]
\begin{align*}
-\pi\frac{m}{A}\,\text{sin}\,\pi\frac{mx}{A}\Bigl[\,&a'\Bigl\{\struck{\text{sin}}\,\pi.\frac{mx}{A} + \frac{3((\Pi-\Pi')+3)}{2.3}\,\text{sin}\,\pi.\frac{mx}{A}\,\text{cos}^2\pi\frac{mx}{A}\\
&\qquad + \frac{5((\Pi-\Pi')+3)((\Pi-\Pi')+3.5))}{2.3.4.5.}\,\text{cos}^5\pi\frac{mx}{A} + \&c\Bigr\}\\
&+ b'\Bigl\{\frac{2\,\struck{\text{sin}}}{\struck{\Pi-\Pi'}}.\pi.\frac{mx}{A}\,\text{cos}.\pi.\frac{mx}{A} + \frac{4((\Pi-\Pi')+8)}{3.4}.\text{sin}.\pi.\frac{mx}{A}\,\text{cos}^3\pi.\frac{mx}{A}\\
&\qquad + \frac{6((\Pi-\Pi')+8)((\Pi-\Pi')+3.8)}{3.4.5.6.}\,\text{sin}\,\pi\frac{mx}{A}\,\text{cos}^5\pi\frac{mx}{A}\Bigr\}\\
&\struck{\times\Bigl[A'\Bigl\{\text{sin}.\pi\frac{ny}{A} + \frac{3((\Pi-\Pi')+3)}{2.3}\,\text{sin}\,\pi\frac{ny}{A}\,\text{cos}^2\pi\frac{ny}{A}}\\
&\qquad \struck{+ \frac{5((\Pi-\Pi')+3)((\Pi-\Pi')+3.5))}{2.3.4.5}\,\text{cos}^5\pi\frac{ny}{A} + \&c\Bigr\}}\\
&\qquad \struck{+ B'\Bigl\{2\,\text{sin}\,\pi\frac{ny}{A}\,\text{cos}\,\pi\frac{ny}{A} + \frac{4((\Pi-\Pi')+8)}{3.4}\,\text{sin}\,\pi.\frac{ny}{A}\,\text{cos}^3\pi\frac{ny}{A}\Bigr\}}\\
&+\,a\,\pi\frac{\sqrt{(m^2+n^2)}}{A},
\end{align*}
\note{La première ligne, « $\frac{dz}{dx} = \ldots\{\ldots\ldots$ », est
laissée en suspens, suivie de points. Dans la ligne en $a'$, « sin »
est écrit sur un autre mot, surchargé ; dans la ligne en $b'$, le
numérateur « sin » et le dénominateur sont noircis d'encre, et sont
donnés comme ratures. La ligne en $A'$, $B'$, tout entière, est barrée
d'un trait ondulé ; sous son deuxième terme est écrit « $4n$ » (lu
« $un$ » peut-être), avec un trait brisé au-dessous. Le « $+$ » isolé
qui suit à droite, sous la ligne barrée, n'est suivi de rien. Les
parenthèses doublées ou non fermées sont celles de la page.}

\begin{align*}
z &= \text{sin}\frac{xm}{A}\ \text{sin}\frac{ym}{A}\left\{\text{cos}\sqrt{(m^2+n^2)}\frac{x}{A} + \text{cos}\sqrt{(m^2+n^2)}\frac{y}{A}\right\}\\
\frac{dz}{dx} &= \frac{m}{A}\,\text{cos}\frac{xm}{A}\ \text{sin}\frac{ym}{A}\left\{\text{cos}\sqrt{m^2+n^2}\frac{x}{A} + \text{cos}\sqrt{m^2+n^2}\frac{y}{A}\right\}\\
&\quad - \frac{\sqrt{m^2+n^2}}{A}\,\text{sin}\frac{xm}{A}\ \struck{\text{cos}}\frac{ym}{A}\left\{\text{sin}\sqrt{(m^2+n^2)}\frac{x}{A}\right\}\\
\frac{ddz}{dx^2} &= -\frac{m^2}{A^2}\,\text{sin}\frac{ym}{A}.\text{sin}\frac{xm}{A}\left\{\text{cos}\sqrt{m^2+n^2}\frac{x}{A} + \text{cos}\sqrt{m^2+n^2}\frac{y}{A}\right\}\\
&\quad - \frac{m^2+n^2}{A}\,\text{sin}\frac{ym}{A}\ \text{sin}\frac{xm}{A}\left\{\struck{\text{cos}}\sqrt{(m^2+n^2)}\frac{x}{A}\right.
\end{align*}
\note{Ces quatre formules, en bas à droite du feuillet, sont d'une
écriture plus fine. Les dénominateurs $A$ de plusieurs fractions sont
barrés d'un trait oblique, comme aux vues précédentes. Dans la troisième
ligne, le premier facteur « sin » est écrit sur un autre mot, et le
second, lu « cos », est surchargé, ce qui est donné comme rature ; de
même le « cos » de la dernière ligne, noirci. La dernière ligne s'arrête
au bord inférieur du feuillet, l'accolade non fermée. Le signe de la
deuxième ligne est un « $=$ » dont le premier trait est plus épais, lu
peut-être « $=-$ ».}
\page{411}

\note{Suite des calculs du feuillet 201 (vue 410), au verso, écrits
cette fois dans le sens de la page ; ni nombre en haut à droite ni
pagination. On y écrit $\text{cos}\,nx$, $\text{sin}\,my$ sans $\pi$ ni
$A$. Plusieurs signes de tête de ligne sont noircis d'encre (un « $-$ »
repassé sur un autre signe) ; ils sont donnés comme on les lit. La
disposition en accolades est simplifiée.}

\begin{align*}
\frac{dz}{dx} = \ill.\,\text{cos}.nx\ \text{sin}\,my\Bigl\{\,&a'\Bigl\{\text{cos}\,nx + \frac{(\Pi-\Pi')+3}{2.3}\,\text{cos}^3nx\\
&\qquad + \frac{((\Pi-\Pi')+3)((\Pi-\Pi')+3.5)}{2.3.4.5}\,\text{cos}^5nx + \&c\Bigr\}\\
&+b'\Bigl\{\frac{2}{\Pi-\Pi'} + \text{cos}^2nx + \&c\Bigr\}\\
-\,n\,\text{sin}\,nx\ \text{sin}\,my\Bigl\{\,&a'\Bigl\{\text{sin}\,nx + \frac{[(\Pi-\Pi')+3]\,3}{2.3}\,\text{cos}^2nx\,\text{sin}\,nx\\
&\qquad + \frac{[((\Pi-\Pi')+3)((\Pi-\Pi')+3.5)]\,5}{2.3.4.5}\,\text{cos}^4nx\,\text{sin}.nx + \&c\Bigr\}\\
&+b'\Bigl\{2\,\text{sin}\,nx\,\text{cos}\,nx + \frac{[(\Pi-\Pi')+8]\,4}{3.4}\,\text{sin}\,nx\,\text{cos}^3nx\\
&\qquad + \frac{[(\Pi-\Pi')+8][(\Pi-\Pi')+3.8]\,6}{3.4.5.6}\,\text{cos}^5nx\,\text{sin}\,nx + \&c
\end{align*}
\note{Le facteur qui suit « $\frac{dz}{dx} =$ » est une lettre noircie,
peut-être $n$ écrit sur un autre signe ; il n'a pas été lu. Le signe
qui ouvre la troisième ligne est un gros trait noirci, lu « $-$ ». Le
premier chiffre du dénominateur « $3.4.5.6$ » est écrit sur un autre.}

\begin{align*}
\frac{ddz}{dx^2} = -\,n^2\,\text{sin}\,nx\ \text{sin}.my\Bigl\{\,&a'\Bigl\{\text{cos}\,nx + \frac{(\Pi-\Pi')+3}{2.3}\,\text{cos}^3nx\\
&\qquad + \frac{((\Pi-\Pi')+3)((\Pi-\Pi')+3.5)}{2.3.4.5}\,\text{cos}^5nx + \&c\Bigr\}\\
&+b'\Bigl\{\frac{2}{\Pi-\Pi'}\ \struck{+}\ \text{cos}^2nx + \frac{(\Pi-\Pi')+8}{3.4}\,\text{cos}^4nx + \&c.\Bigr\}\\
-\,\struck{2}n^2\,\text{sin}\,nx\ \text{sin}\,my\Bigl\{\,&a'\Bigl\{\text{cos}\,nx + \frac{[.(\Pi-\Pi')+3]\,3}{2.3}.\text{cos}^3nx\\
&\qquad \struck{+}\ \frac{[((\Pi-\Pi')+3)((\Pi-\Pi')+3.5)]\,5}{2.3.4.5}\,\text{cos}^5nx + \&c\\
&\qquad - [(\Pi-\Pi')+3]\,\text{cos}\,nx\,\text{sin}^2nx\\
&\qquad - \frac{((\Pi-\Pi')+3)((\Pi-\Pi')+3.5)}{2.3}\,\text{cos}^3nx\,\text{sin}^2 - \&c.\Bigr\}\\
&+b'\Bigl\{2(\text{cos}^2nx - \text{sin}^2nx) + \frac{[(\Pi-\Pi')+8]}{3.4}.4\,\text{cos}^4nx\\
&\qquad + \frac{[(\Pi-\Pi')+8][(\Pi-\Pi')+3.8]}{3.4.5.6}\,6\,\text{cos}^6nx + \&c.\\
&\qquad \struck{+}\,[(\Pi-\Pi')+8]\,\text{sin}^2nx\,\text{cos}^2nx\\
&\qquad - \frac{[(\Pi-\Pi')+8][(\Pi-\Pi')+3.8]}{3.4.}\,\text{sin}^2nx\,\text{cos}^4nx - \&c
\end{align*}
\note{Après « $\frac{2}{\Pi-\Pi'}$ », à la deuxième accolade, un « $+$ »
est noirci ; de même le « $+$ » devant le terme en $\text{cos}^5nx$ de
la ligne suivante et celui qui ouvre la dernière ligne. Le 2 devant
« $n^2$ », au début du second bloc, est écrit sur un autre caractère
et barré ; il reparaît, net, au bloc suivant. Le dernier facteur de la
ligne en $\text{cos}^3nx\,\text{sin}^2$ est coupé au bord droit.}

\begin{align*}
-\,2n^2\,\text{sin}\,nx\ \text{sin}\,my\Bigl\{\,&a'\Bigl\{\text{cos}\,nx + \frac{[(\Pi-\Pi')+3]}{2.3}\,3\,\text{cos}^3nx\\
&\qquad + \frac{[((\Pi-\Pi')+3)((\Pi-\Pi')+3.5)].5}{2.3.4.5}\,\text{cos}^5nx + \&c\\
&+b'\Bigl\{2\,\text{cos}^2nx + \frac{[(\Pi-\Pi')+8]}{3.4}\,4\,\text{cos}^4nx\\
&\qquad + \frac{[(\Pi-\Pi')+8][(\Pi-\Pi')+3.8]}{3.4.5.6}\,6.\text{cos}^6nx + \&c\Bigr\}
\end{align*}
\begin{align*}
= n^2\,\text{sin}\,nx\ \text{sin}.my\Bigl\{\,&a'\Bigl\{((\Pi-\Pi')-1)\,\text{cos}\,nx\\
&\qquad + \frac{((\Pi-\Pi')+3)}{2.3}\,[(\Pi-\Pi')+3.5-1-9-6]\,\text{cos}^3nx\\
&\qquad + \frac{((\Pi-\Pi')+3)((\Pi-\Pi')+3.5)}{2.3.4.5}\\
&\qquad\quad \times[(\Pi-\Pi')+7.5-1-25-10]\,\text{cos}^5nx + \&c.\Bigr\}\\
&+b'\Bigl\{\frac{2}{\Pi-\Pi'}\ \struck{+}\ 2 + ((\Pi-\Pi')+8)\ \struck{-}\ 4-4-1)\,\text{cos}^2nx\\
&\qquad + \frac{((\Pi-\Pi')+8)}{3.4.}\,((\Pi-\Pi')+3.8\ \struck{+}\ 4\ \struck{+}\ 12\ \struck{+}\ 8-1)\,\text{cos}^4nx + \&c
\end{align*}
\note{Dans la dernière égalité, le signe qui suit l'accolade, devant
$a'$, est une grosse tache d'encre, et un autre trait épais tombe après
« $\text{cos}\,nx$ » ; ils ne sont pas repris. Les signes entre les
nombres de la ligne en $b'$ (« $\frac{2}{\Pi-\Pi'}\ +\ 2$ », « $-\ 4-4-1$ »)
et de la dernière ligne (« $+4+12+8-1$ ») sont des « $+$ » noircis ou
changés en « $-$ » ; les « $+$ » sont donnés comme ratures, et le sens de
chaque correction n'est pas sûr. « $+7.5$ » : le 7 est repassé sur un
autre chiffre. Le bas du feuillet est blanc.}
\page{412}

ainsi, lorsque l'on fera $z' = \ldots\text{sin}\,\pi.\frac{xm}{A}\ \text{sin}\,\pi.\frac{ym}{A}$, on pourra prendre
\begin{align*}
z = \ldots\ \text{sin}\,\pi.\frac{xm}{A}\ \text{sin}\,\pi.\frac{ym}{A}\Bigl(\\
&a'\Bigl\{\text{cos}.\pi.\frac{mx}{A} + \frac{\Pi+3}{2.3}\,\text{cos}^3\pi.\frac{mx}{A} + \ldots\ldots\\
&\quad + \frac{(\Pi+3)(\Pi+3.5)\ldots(\Pi+4n^2-1)}{2.3.4.5\ldots 2n+1}\,\text{cos}^{2n+1}\pi.\frac{mx}{A} + \&c\Bigr\}\\
&+b'\Bigl\{\frac{2}{\Pi} + \text{cos}^2\pi.\frac{mx}{A} + \frac{\Pi+8}{3.4}\,\text{cos}^4\pi.\frac{mx}{A}\ldots\\
&\quad + \frac{(\Pi+8)(\Pi+3.8)\ldots(\Pi+(2n-1)^2-1)}{3.4.5.6\ldots 2n}\,\text{cos}^{2n}\pi.\frac{mx}{A} + \&c\Bigr\}\\
\times\ &A'\Bigl\{\text{cos}.\pi.\frac{my}{A} + \frac{\Pi'+3}{2.3}\,\text{cos}^3\pi.\frac{my}{A} + \ldots\\
&\quad + \frac{(\Pi'+3)(\Pi'+3.5)\ldots(\Pi'+4n^2-1)}{2.3.4.5\ldots 2n+1}\,\text{cos}^{2n+1}\pi.\frac{my}{A} + \&c\Bigr\}\\
&+B'\Bigl\{\frac{2}{\Pi'} + \text{cos}^2\pi.\frac{my}{A} + \frac{\Pi'+8}{3.4}.\text{cos}^4\pi.\frac{my}{A}\ \struck{\ldots}\\
&\quad + \frac{(\Pi'+8)(\Pi'+3.8)\ldots(\Pi'+(2n-1)^2-1)}{3.4.5.6\ldots 2n}\,\text{cos}^{2n}\pi.\frac{my}{A} + \&c\Bigr\}\\
&+a''\Bigl\{\text{cos}\,\pi.\frac{mx}{A} + \frac{\Pi''+3}{2.3}\,\text{cos}^3\pi.\frac{mx}{A} + \ldots\\
&\quad + \frac{(\Pi''+3)(\Pi''+3.5)\ldots(\Pi''+4n^2-1)}{2.3.4.5\ldots 2n+1}.\text{cos}^{2n+1}\pi.\frac{mx}{A} + \&c\Bigr\}\\
&+b''\Bigl\{\frac{2}{\Pi''} + \text{cos}^2\pi.\frac{mx}{A} + \frac{\Pi''+8}{3.4}.\text{cos}^4\pi.\frac{mx}{A} + \ldots\\
&\quad + \frac{(\Pi''+8)(\Pi''+3.8)\ldots(\Pi''+(2n-1)^2-1)}{2.3.4.5\ldots 2n}\,\text{cos}^{2n}\pi.\frac{mx}{A} + \&c\Bigr\}\\
&+A''\Bigl\{\text{cos}\,\pi.\frac{my}{A} + \frac{\Pi''+3}{2.3}\,\text{cos}^3\pi.\frac{my}{A} + \ldots\\
&\quad + \frac{(\Pi''+3)(\Pi''+3.5)\ldots(\Pi''+4n^2-1)}{2.3.4.5\ldots 2n+1}\,\text{cos}^{2n+1}\pi.\frac{my}{A} + \&c\Bigr\}\\
&+B''\Bigl\{\frac{2}{\Pi''} + \text{cos}^2\pi.\frac{my}{A} + \frac{\Pi''+8}{3.4}\,\text{cos}^4\pi.\frac{my}{A} + \ldots\\
&\quad + \frac{(\Pi''+8)(\Pi''+3.8)\ldots(\Pi''+(2n-1)^2-1)}{2.3.4.5\ldots 2n}\,\text{cos}^{2n}\pi.\frac{my}{A} + \&c\Bigr\}\\
&\Bigr)\ \ldots\ (G)
\end{align*}
\note{En tête, vers le milieu, « (11) », le chiffre barré d'un trait ;
en haut à droite « 202 », à l'encre, de la série courante. Ce feuillet
est une nouvelle mise au net de la fin de la page (9) de la vue 408 et
du début de la page (10) de la vue 409 : la même formule (G), écrite
« $\pi.\frac{mx}{A}$ » sans parenthèses, puis la même phrase et le début
de (H). Deux différences avec la vue 408 : les dénominateurs des lignes
en $b''$ et $B''$ se lisent « $2.3.4.5\ldots 2n$ » et non
« $3.4.5.6\ldots 2n$ » ; après le troisième terme de la ligne en $B'$,
une ligne de points est noircie. Au-dessus du « (11) », en haut à gauche
du chiffre 202, une figure pâle entre parenthèses, « (12) » peut-être,
est sans doute celle du verso vue par transparence. Disposition de la
formule simplifiée, comme aux vues 407 et 408.}

pourvu que $\Pi$ $\Pi'$ et $\Pi''$ \struck{\uncertain{aient}} \add{ayent} entre elles la relation exprimée
par l'équation $\Pi+\Pi' = \struck{\ill}\,\Pi''$.

Les sons simples seront comme dans les \struck{valeurs} cas compris
dans l'intégrale (F), c'est adire que l'on aura toujours
$\sqrt{\tfrac{1}{K}}.\frac{1}{\pi} = \frac{\pi.m^2(\Pi+\Pi'-2)}{A^2}$.

\struck{Et} lorsque \add{l'on} prendra $z' = \text{sin}.\pi.\frac{nx}{A}\ \text{sin}.\pi.\frac{my}{A}$, on pourra faire
\begin{align*}
z = \ldots\ \text{sin}.\pi.\frac{xn}{A}\ \text{sin}.\pi.\frac{ym}{A}\Bigl(\,&a'\Bigl\{\text{cos}.\pi.\frac{nx}{A} + \frac{2-m^2+3}{2.3}\,\text{cos}^3\pi.\frac{nx}{A}\\
&\qquad + \frac{(2-m^2+3)(2-m^2+15)}{2.3.4.5}\,\text{cos}^5\pi.\frac{nx}{A} + \&c\Bigr\}\\
&+b'\Bigl\{\frac{2}{2-m^2} + \text{cos}^2\pi.\frac{nx}{A} + \frac{2-m^2+8}{3.4}\,\text{cos}^4\pi.\frac{nx}{A}\\
&\qquad + \frac{(2-m^2+8)(2-m^2+24)}{3.4.5.6}\,\text{cos}^6\pi.\frac{nx}{A} + \&c\Bigr\}\\
&+A'\Bigl\{\text{cos}\,\pi.\frac{my}{A} + \frac{2-n^2+3}{2.3}.\text{cos}^3\pi.\frac{my}{A}\\
&\qquad + \frac{(2-n^2+3)(2-n^2+15)}{2.3.4.5.}\,\text{cos}^5\pi.\frac{my}{A} + \&c\Bigr\}\\
&+B'\Bigl\{\frac{2}{2-n^2} + \text{cos}^2\pi.\frac{my}{A} + \frac{2-n^2+8}{3.4}\,\text{cos}^4\pi.\frac{my}{A}\\
&\qquad + \frac{(2-n^2+8)(2-n^2+\uncertain{24})}{3.4.5.6}\,\text{cos}^6\pi.\frac{my}{A} + \&c\Bigr\}\Bigr)\ \ldots\ (H)
\end{align*}
\note{« ayent » est écrit au-dessus d'un mot barré, lu « aient » avec doute. Après « $\Pi+\Pi' =$ », un groupe (un 2 et une lettre, semble-t-il)
est noyé sous une rature en boucle, et « $\Pi''$ » est écrit au-dessus ;
ce que la rature couvre n'est pas lu. « valeurs » est barré devant
« cas ». « Et » est barré en tête de la phrase et « l'on » ajouté
au-dessus de « prendra ». Dans (H), les constantes « $+3$ », « $+8$ »,
« $+15$ », « $+24$ » sont presque toutes repassées d'une encre épaisse
sur d'autres chiffres, et le premier chiffre du dénominateur de
« $\frac{2}{2-m^2}$ » est surchargé ; une grosse tache d'encre couvre la
fin du numérateur du dernier terme de la ligne en $B'$, lu « $+24$ »
d'après les lignes voisines. Les textes de la vue 409 écrivent ces
coefficients avec une parenthèse ouverte, « $(2-m^2+3$ ». La formule
occupe le bas du feuillet, qui s'arrête là.}
\page{413}

Les sons simples qui correspondent aux manières de vibrer renfermées
dans cette dernière intégrale, seront exprimés par la formule
$\sqrt{\tfrac{1}{K}}.\frac{1}{\pi} = \frac{\pi.m^2n^2\tau\sqrt{b}}{A^2}$.
\note{En tête, au milieu, « (12) » ; aucun nombre en haut à droite de ce
côté du feuillet, verso du feuillet 202 (vue 412). À côté du « (12) », un
« (11) » pâle et retourné est celui du recto, vu par transparence. Cette
page est une nouvelle mise au net de la seconde moitié de la page (10)
de la vue 409, avec quelques différences notées plus bas. Le « $\tau$ »
de l'épaisseur est tracé ici nettement comme un $\Gamma$.}

N.o 6 En prenant encore $z = z'z''$, et fesant
$z' = \text{sin}(t.\sqrt{\tfrac{1}{K}}+\zeta)\ \text{sin}\,\pi.\frac{nx}{A}$, on aura \struck{encore} toujours une
équation dela forme de celle aux cordes vibrantes
savoir $\left(\frac{ddz'}{dt^2}\right) = \left(\frac{ddz'}{dx^2}\right)\frac{K}{n^2}$. Elle indique que la surface
sera periodique dans le sens de $x$. On pourra
faire
\begin{align*}
z = \ldots\ \text{sin}.\pi.\frac{nx}{A}\Bigl\{\,&a'\Bigl\{\text{cos}.\pi.\frac{nx}{A} + \frac{2-m^2+3}{2.3.}\,\text{cos}^3\pi.\frac{nx}{A}\\
&\qquad + \frac{(2-m^2+3)(2-m^2+15)}{2.3.4.5}\,\text{cos}^5\pi.\frac{nx}{A} + \&c\Bigr\}\\
&b'\Bigl\{\frac{2}{-m^2+2} + \text{cos}^2\pi.\frac{nx}{A} + \frac{2-m^2+8}{3.4}\,\text{cos}^4\pi.\frac{nx}{A} + \&c\Bigr\}\\
&+\ \mathfrak{A}\,\text{sin}.\pi.\frac{\uncertain{mn}.y}{A}\ \ldots\ldots\ (I)
\end{align*}
\note{L'étiquette « (I) » est écrite sur une autre lettre, noircie. Les
constantes « $+3$ », « $+8$ » et « $+15$ » sont repassées sur d'autres
chiffres, comme à la vue 412. L'argument du dernier sinus se lit
« $\pi.mn.y$ » avec doute (à la vue 409, « $\pi\,may$ ») ; la majuscule
ornée est la même.}

\struck{Les diff} mais cette \add{surface} ne jouira pas dela même propriété
dans le sens des $y$ puisque cette dernière variable n'entre
pas dans l'équation aux cordes vibrantes.
Les sons simples seront exprimés par la formule $\sqrt{\tfrac{1}{K}}.\frac{1}{\pi} = \frac{\pi m^2n^2\tau\sqrt{b}}{A^2}$
\note{« Les diff » est barré en tête de la phrase, qui reprend « mais
cette… » ; « surface » est ajouté au-dessus de « ne jouira », écrit sur un
autre mot. La ligne sur les sons simples est serrée, d'une écriture plus
petite, entre deux lignes ; le dénominateur « $A^2$ » y est surchargé. La
vue 409 barrait cette même valeur pour écrire « $= \pi.n^2(m^2+1)$ » ;
ici elle est gardée.}

Les differentes intégrales que je viens de donner, conviennent
toutes, excepté la première, à l'équation differentielle des
surfaces tendues; et c'est a quoi elles doivent leur simplicité.
car si on vouloit avoir les intégrales qui n'appartiennent
qu'à l'équation differentielle des surfaces élastiques, \struck{on}
\struck{\uncertain{faudroit} avoir} \add{on trouveroit} nonseulement \struck{les} termes contenant $e^{\pm\sqrt{-1}.\omega}$,
mais \struck{\uncertain{sur}} encore des termes où entreroient, de la même
manière, les \struck{\ill{}} \add{exponentielles} quantités $e^{\pm\omega}$; et alors l'angle
$\omega$ ne pourroit être déterminé que par des calculs
d'aproximation; desorte que non seulement les
\note{Au-dessus de deux mots barrés au début de la ligne (le premier lu
« faudroit » avec doute, le second « avoir »), elle écrit « on
trouveroit » ; « on » est barré en fin de ligne précédente. Le « les »
barré devant « termes » est noyé d'encre. Le mot barré sous
« exponentielles », ajouté au-dessus de la ligne, semble une première
graphie du même mot, trop raturée pour être lue ; « quantités » est gardé sur la ligne. Le mot barré
après « mais » est court et n'est lu qu'avec doute. La page s'arrête sur
« les », suivi d'un trait qui appelle la suite.}
\page{414}

intégrales seroient beaucoup plus compliquées que les précédentes;
mais encore elles seroient d'une application moins facile que celles-ci.
Avant de m'occuper dela comparaison entre les résultats de la
théorie \struck{et celle} del'experience, je vais présenter quelques remarques
que je crois importantes.
\note{En tête, vers le milieu, « (13) », le chiffre barré d'un trait,
comme le « (11) » de la vue 412 ; en haut à droite « 203 », à l'encre,
de la série courante. La première ligne achève la phrase laissée sur
« les » au bas de la vue 413 (« desorte que non seulement les intégrales
seroient… »). « et celle » est barré, « de » étant rattaché à
« l'experience ». Un trait vertical court, dans la marge de droite, face
à « remarques », n'appelle aucune note.}

\section{§. 3. Observations sur les propriétés analytiques des differens points des corps sonores}

N.o 7. Indépendamment des conditions auxquelles doivent satisfaire
tous les points renfermés dans l'étendue des intégrales,
il en existe d'autres auxquelles les limites de ces intégrales
sont assujetties, et qui parconséquent appartiennent \struck{à tous} \add{à} les
points du contour des corps sonores, à l'exclusion des autres
points des mêmes corps. Il suit delà qu'aucun point ou ligne
ne pourront \struck{dans aucun cas} \add{jamais} être considérés comme limite
d'une partie vibrante, \struck{si} s'ils ne satisfont aux conditions
analytiques qui distinguent les points du contour des autres
points du corps sonore; et qu'au contraire il suffira toujours
qu'un point ou ligne satisfassent a ces mêmes conditions,
pour qu'il soit permis de les considérer comme limites
analytiques des parties vibrantes qu'ils séparent.
\note{Le titre du § 3 est souligné au milieu d'un trait court ; son
dernier mot, « sonores », est serré contre la marge. Dans « appartiennent
à tous les », « à tous » est surchargé plutôt que barré : un « à » est
écrit sur « tous », et la lecture de la correction reste douteuse.
« jamais » est écrit au-dessus de « dans aucun cas », dont « aucun cas »
seul est nettement barré ; « dans » n'est peut-être pas annulé. « si »
est écrit sur un autre mot devant « s'ils ».}

Lorsqu'il s'agit dela corde vibrante, on sait depuis longtems
que tous \add{les} points de repos \struck{peuvent} \add{peut} être considérés comme \add{de veritables} limites du
corps sonore: mais ces points ne sont pas les seuls qui
satisfassent aux conditions des extrémités, et il est aisé de voir
que tous les points dans lesquels la tangente est parallèle à l'axe
\note{Correction malaisée à suivre : « les » est ajouté au-dessus de
« tous », le « s » de « points » semble repassé, et « peuvent » est
corrigé à la main en un mot plus court, lu « peut » avec doute ; « de
veritables » est ajouté au-dessus de « limites ». Le texte s'arrête sur
« à l'axe », au bas du feuillet, dont le dernier tiers est un carton de
montage blanc. La phrase se poursuit en tête de la vue 415.}
\page{415}

c'est-a-dire tous les points situés au milieu des ventres de
vibration, jouissent dela même propriété. En effet, d'après
les suppositions admises dans la théorie du son, on peut conclure
de la formule P 332 de la 1\textsuperscript{ere} édition de la Mec. Analy.
que les conditions des limites se réduisent à l'équation
$R.\frac{dz}{dx}.\delta z = 0$. àlaquelle on peut également satisfaire par
l'une ou l'autre des équations $z = 0$, $\frac{dz}{dx} = 0$. A la verité
il y a une impossibilité physique à vérifier ce résultat pour
la corde tendue; mais il reçoit une parfaite confirmation
dans la théorie des instrumens à vent; car on peut voir dans
le second volume des mémoires de Turin que M.\textsuperscript{r} Dela
Grange a conclu l'équation $\frac{dz}{dx} = 0$ dela nécessité
que l'élasticité de l'air dans le tuyau soit la même
à l'extrémité ouverte que celle de l'atmosphere. Ainsi,
tandis que les extrémités dela corde tendue fournissent une
manière desatisfaire àla condition $R.\frac{dz}{dx}.\delta z = 0$, l'autre est
donné par les extrémités dela flûte ouverte aux deux bouts.
desorte qu'en prenant de part et d'autre le cas le
plus simple, l'un est exactement le renversement de
l'autre; car lorsque la corde vibre dans toute son étendue
en ne formant qu'une seule période, il existe trois points
\struck{d'extrémités} \struck{\uncertain{de}} limites savoir: les deux points d'extrémités
effectives \struck{lesquels} pour lesquels on a $z = 0$ et le point de
milieu pour lequel $\frac{dz}{dx} = 0$, tandis que pour la manière
\note{En tête, au milieu, « (14) », le 4 écrit sur un autre chiffre ;
aucun nombre en haut à droite de ce côté du feuillet, verso du
feuillet 203 (vue 414). La première ligne achève la phrase de la vue
414 (« tous les points dans lesquels la tangente est parallèle à l'axe
c'est-a-dire tous les points situés au milieu des ventres… »). « P 332 »
est écrit ainsi, le P majuscule pour « page » ; « Mec. Analy. » est
l'abréviation de la page, pour la Mécanique analytique. Le signe de
variation, dans « $R.\frac{dz}{dx}.\delta z$ », est la lettre bouclée
rendue par $\delta$ ; la première fois il est écrit sur un autre signe.
« Dela Grange » est écrit en deux mots, coupé en fin de ligne. Devant
« limites », « d'extrémités » est barré, suivi d'une ou deux lettres
barrées et noircies ; « lesquels » est barré avant « pour lesquels ». Le
texte s'arrête sur « la manière » au bas du feuillet, le dernier tiers
de la vue étant le carton de montage. La suite est à la vue 418, le
petit feuillet de la vue 416 (et son verso blanc, vue 417) s'intercalant
entre les deux.}
\page{416}

Il regne dans tous le passage relatif
aux effets de l'application dela cire
une erreur qui m'a d'abord inquiété
lorsqu'on a voulu aprecier le son dela
plaque chargée en supposant le
poids de cette dernière a celui dela
plaque non chargée comme $N$ a $n$ on a
pris le son proportionnel au nombre
$\frac{N^2+n^2}{n^2}$ parconséquent ce son seroit $>$ aigu que
$\frac{2n^2}{n^2}$ qui est celui dela plaque non chargée
le contraire doit avoir lieu et le premier son
est proportionnel au nombre $\frac{N^2+n^2}{N^2}$ tandis que
$\frac{N^2+n^2}{n^2}$ exprimeroit le
son d'une plaque de
$N$ grandeur dont on auroit
retranchée $(N-n)$
La difference que j'ai
exprimé par $\frac{N^2+n^2}{2n^2}$
doit l'être par
$\frac{2N^2}{N^2+n^2}$ qui heureusement
ne different pas sensiblement
l'un de l'autre.
\note{Une pièce nouvelle, sur un petit feuillet déchiré, collé sur un
carton de montage : en haut à droite « 204 », à l'encre, de la série
courante ; aucune pagination de sa main. L'écriture est d'une plume plus
fine et d'une encre plus pâle, rapide, sans rature. Le texte occupe
d'abord toute la largeur du feuillet, puis, à partir de « exprimeroit
le », une colonne étroite sur la droite d'un prolongement du papier ;
c'est une note isolée sur une expérience de plaque chargée de cire, sans
lien visible avec les pages (4) à (14) qui précèdent, et qui ne reprend
pas la phrase laissée sur « la manière » à la vue 415. « $>$ aigu » est
écrit ainsi, le signe tenant lieu de « plus ». « le son » est écrit sur
un autre mot, noirci ; « a » pour « à » deux fois (« a celui »,
« comme $N$ a $n$ »). Le « $2$ » de « $\frac{2n^2}{n^2}$ » et de
« $\frac{N^2+n^2}{2n^2}$ » est tracé à la façon d'un $\mathcal{B}$, comme
ailleurs dans ce volume.}
\page{418}

de vibrer la plus simple dont soit susceptible une flûte ouverte
aux deux bouts, les extrémités effectives satisfont à la condition
$\frac{dz}{dx} = 0$, et le point de milieu auquel se reunissent les deux
demi periodes, remplit la condition $z = 0$.
\note{En tête, à gauche du milieu, « (15) » ; en haut à droite « 205 », à
l'encre, de la série courante. La première ligne achève la phrase
laissée au bas de la vue 415 (« tandis que pour la manière de vibrer la
plus simple… ») : le petit feuillet de la vue 416, avec son verso blanc
(vue 417), est intercalé entre les pages (14) et (15).}

Quelque chose de semblable a lieu pour les surfaces
tendues; et c'est cequ'on peut voir aisement en introduisant
\struck{les simplifications admises}
dans les formules données par M.\textsuperscript{r} Biot P. 87 deson memoire
t. 4. del'institut, les simplifications admises dans la théorie du
son: car les conditions des extrémités se réduisent alors
à $\int E.dx.\frac{dz}{dy}.\delta z = 0$ et $\int E.dy.\frac{dz}{dx}\,\delta z = 0$. \struck{et} \add{Or,} dans le cas
des surfaces rectangulaires, un des côtés peut toujours etre
supposé parallèle à l'axe des $x$ ou a celui des $y$;
desorte qu'une des conditions cidessus est nécessairement toujours
satisfaite par l'une des deux suppositions $dx = 0$, $dy = 0$,
tandis que l'autre doit l'être par l'une des conditions
$z = 0$, ou $\frac{dz}{dy}$; $z = 0$ ou $\frac{dz}{dx} = 0$. l'équation $z = 0$ appartient
aux extrémités effectives ou aux lignes de repos dont M\textsuperscript{r}
Biot a demontré l'existence; et les equations $\frac{dz}{dy} = 0$ $\frac{dz}{dx} = 0$
montrent qu'il n'est pas toujours nécessaire de faire
$E = 0$, pour satisfaire aux conditions des limites dans des
points qui ne sont pas supposés fixes.
\note{« les simplifications admises » est barré au début d'une ligne, puis
récrit une ligne plus bas, après la référence. « P. 87 » : le P
majuscule pour « page » ; « t. 4. » pour le tome 4 des mémoires de
l'Institut. Les deux intégrales sont écrites par un s long, rendu par
$\int$ ; le signe de variation est la lettre bouclée rendue par
$\delta$. « Or, » est écrit au-dessus de « et », barré. Dans
« nécessairement », les premières lettres sont écrites sur d'autres,
noircies ; la lecture est sûre par le sens. « ou $\frac{dz}{dy}$ » est
écrit sans « $= 0$ ».}

On voit donc que, quelque soit le corps sonore que l'on
considère, les points ou les lignes de repos ne sont pas
les seuls \struck{points} ou lignes qui satisfassent aux conditions
des extrémités, et qui soient parconséquent de veritables
\note{« points » est barré devant « ou lignes ». Le texte s'arrête sur
« de veritables » au bas du feuillet ; la suite est à la vue 419.}
\page{419}

limites analytiques des parties vibrantes qu'ils séparent.
\note{En tête, à gauche du milieu, « (16) » ; aucun nombre en haut à
droite de ce côté du feuillet, verso du feuillet 205 (vue 418). La
première ligne achève la phrase de la vue 418 (« …et qui soient
parconséquent de veritables limites analytiques… »).}

On voit d'ailleurs que, tant qu'il ne s'agit que des surfaces
tendues, il est toujours permis de considerer, \struck{les \uncertain{cordes}} ainsi que
l'indique le programme publié par la classe, les \emph{nappes
vibrantes} comme analogues aux \emph{ondes} des cordes de Sauveur,
et les courbes \emph{d'équilibre ou de repos} comme correspondant
aux nœuds ou \emph{points de repos} des mêmes cordes, \struck{car} \add{puisque} de
part et d'autre l'équation $z = 0$ suffit pour remplir
les conditions des extrémités. Il n'en est pas demême
pour les lames et les surfaces élastiques; car il existe
dans leurs mouvemens deux genres differens de points ou
lignes de repos; les uns qui, parcequ'ils satisfont aux
conditions des extrémités, sont delamême nature que ceux
qui se forment sur les cordes et les surfaces tendues,
les autres qui, par la raison contraire, séparent des
espaces qui, \struck{\uncertain{s}} s'ils étoient isolés, ne pourroient, dans
aucun cas, se ployer aux courbures qu'ils prennent
comme parties integrantes du corps vibrant ni parconséquent
rendre le même son que lui.
\note{Les mots donnés en italique sont soulignés sur la page : « nappes
vibrantes », « ondes », « d'équilibre ou de repos », « points de
repos ». Après « considerer, », deux mots sont barrés, le second lu
« cordes » avec doute. « puisque » est écrit au-dessus de « car », barré.
« son », dans « le même son », est écrit sur un autre mot, noirci. « la
classe » : la Classe des sciences de l'Institut.}

Quoique la remarque que je viens d'exposer n'ait pas
encore été faite, elle se présente naturellement en
examinant la théorie qu'Euler a donnée pour la lame
élastique \emph{Acta academiæ} $\ldots$ \emph{Petropolitanæ} pro anno 1779. Pars prior:
c'est cequi je vais établir dans le N.o suivant.
\note{Le nom d'Euler est tracé ici comme « Lobr » : c'est son E initial
en forme de L, déjà relevé aux vues 210 à 218, et la fin du nom, très
cursive ; la lecture « Euler » est assurée par la référence qui suit.
« Acta academiæ » est souligné, suivi d'une ligne de points soulignée
et de « Petropolitanæ », souligné, dont les lettres du milieu sont
malaisées ; « pro anno 1779. Pars prior » est d'une écriture plus
petite. « cequi » est écrit ainsi pour « ce que ». Le bas du feuillet
est occupé par le carton de montage, sur lequel se lisent à l'envers,
par transparence, quelques mots du petit feuillet de la vue 416 (« doit
l'être par », « $2N^2$ », « qui heureusement »).}
\page{420}

N.o 8. Pour y parvenir, en adoptant les simplifications admises par Euler,
je prens les conditions des extremités données par M.\textsuperscript{r} Dela Grange
P. \uncertain{197} dela nouvelle éd. de sa mec. Elles se réduisent dans cette hypothèse
aux deux équations $K.\frac{d^3z}{dx^3}.\delta z = 0$ et $K\frac{d^2z}{dx^2}.\delta\frac{dz}{dx} = 0$. je remarque que
l'on doit avoir pour y satisfaire l'un ou l'autre des quatre
couples suivans d'équations. $z = 0$ et $\frac{dz}{dx} = 0$; $z = 0$ et $\frac{d^2z}{dx^2} = 0$;
$\frac{d^3z}{dx^3} = 0$ et $\frac{dz}{dx} = 0$; $\frac{d^3z}{dx^3} = 0$ et $\frac{d^2z}{dx^2} = 0$. et je conclus déja que
la condition $z = 0$ ne suffit pas pour que le point auquel
elle appartient puisse être considéré comme une veritable limite.
\note{En tête, à gauche du milieu, « (17) » ; en haut à droite « 206 », à
l'encre, de la série courante. Ce N.o 8 est celui qu'annonce la fin de
la vue 419 (« c'est cequi je vais établir dans le N.o suivant »). Le nom
d'Euler, en fin de première ligne, est écrit d'une encre épaisse, sur
un autre mot semble-t-il ; la lecture est celle du sens. Le numéro de
page de Lagrange est chargé d'encre : 197 avec doute, le 9 pouvant être
un autre chiffre. « mec. » : la Mécanique analytique ; « la nouvelle
éd. » s'oppose à la « 1\textsuperscript{ere} édition » citée à la vue
415. La lettre devant les fractions des deux équations est un $K$, lu
avec doute pour la seconde. Le signe de variation est la lettre bouclée
rendue par $\delta$. Dans les exposants de $d^3z$, $d^2z$, le chiffre
est posé au-dessus du $d$.}

On peut voir P. 116 du mémoire cité qu'Euler n'a examiné que
trois de ces couples, l'un $z = 0$ et $\frac{dz}{dx}$; qu'il attribue aux
extrémités fixées, l'autre $z = 0$ et $\frac{d^2z}{dx^2}$; qu'\struck{il regarde comme}
qui appartient aux extrémités appuyées, et le troisième
$\frac{d^3z}{dx^3} = 0$ et $\frac{d^2z}{dx^2} = 0$ qu'il regarde comme formant le caractère
des extrémités libres; cependant il est clair que les extrémités
seroient également libres lorsqu'elles appartiendroient au quatrième
couple $\frac{d^3z}{dx^3} = 0$ et $\frac{dz}{dx} = 0$. A la verité, lorsque l'on fait
vibrer une lame physiquement libre à ses extrémités, on obtient bien
plus aisément les figures et les sons qui se rapportent aux
conditions $\frac{d^3z}{dx^3} = 0$ et $\frac{d^2z}{dx^2} = 0$ que ceux qui conviennent aux
conditions $\frac{d^3z}{dx^3} = 0$ et $\frac{dz}{dx} = 0$. Cependant je crois être parvenue
à obtenir ces derniers. mais la lame qui se prêtoit plus
aisément que les autres à cette manière de vibrer, s'est
brisée trop tôt pour que je pusse \struck{\uncertain{mesurer}} constater et mesurer
ce résultat.
\note{« P. 116 du mémoire cité » : le mémoire d'Euler nommé à la vue 419
(Acta Petrop. 1779, pars prior). Le nom est tracé comme à la vue 419, E
en forme de L. « examiné » est écrit sur un autre mot, noirci au début.
« $\frac{dz}{dx}$ » et « $\frac{d^2z}{dx^2}$ », dans l'énumération, sont
écrits sans « $= 0$ ». « il regarde comme » est barré et remplacé, sur
la ligne suivante, par « qui appartient ». Le mot barré après « pusse »
est noirci ; il est lu « mesurer » avec doute, le même mot revenant
après « constater et ». « parvenue » est au féminin sur la page. Le
bas du feuillet est blanc ; la suite n'est pas dans ce lot.}

\end{document}
